%%% LaTeX-Vorlage Version 1.8 %%%

% Grundlegende Dokumenteneigenschaften gemäß DHBW-Vorgaben
\documentclass[a4paper,fontsize=11pt,oneside,parskip=half,headings=normal]{scrreprt} 
% \usepackage{showframe} % nur für Kontrolle der Ränder 

%%% Präambel einbinden (mit Festlegungen gemäß DHBW-Vorgaben) %%%
%%% Präambel %%%
% hier sollten keine Änderungen erforderlich sein
%
\usepackage[utf8]{inputenc}   % Zeichencodierung UTF-8 für Eingabe-Dateien
\usepackage[T1]{fontenc}      % Darstellung von Umlauten im PDF

\usepackage{listings}         % für Einbindung von Code-Listings
\lstset{numbers=left,numberstyle=\tiny,numbersep=5pt,texcl=true}
\lstset{literate=             % erlaubt Sonderzeichen in Code-Listings 
{Ö}{{\"O}}1
{Ä}{{\"A}}1
{Ü}{{\"U}}1
{ß}{{\ss}}2
{ü}{{\"u}}1
{ä}{{\"a}}1
{ö}{{\"o}}1
{€}{{\euro}}1
}

\usepackage[
  inner=35mm,outer=15mm,top=25mm,
  bottom=20mm,foot=12mm,includefoot
]{geometry}                 % Einstellungen für Ränder

\usepackage[ngerman]{babel} % Spracheinstellungen Deutsch
\usepackage[babel,german=quotes]{csquotes} % deutsche Anf.zeichen
\usepackage{enumerate}      % anpassbare Nummerier./Aufz.
\usepackage{graphicx}       % Einbinden von Grafiken
\usepackage[onehalfspacing]{setspace} % anderthalbzeilig

\usepackage{blindtext}      % Textgenerierung für Testzwecke
\usepackage{color}          % Verwendung von Farbe 

\usepackage{acronym}        % für ein Abkürzungsverzeichnis

\usepackage[                % Biblatex
  backend=biber,
  bibstyle=_dhbw_authoryear,maxbibnames=99,
  citestyle=authoryear,     
  uniquename=true, useprefix=true,
  bibencoding=utf8]{biblatex}
%kein Punkt am Ende bei \footcite
%http://www.golatex.de/footcite-ohne-punkt-am-schluss-t4865.html
\renewcommand{\bibfootnotewrapper}[1]{\bibsentence#1}


%Reihenfolge der Autorennamen
%   
% http://golatex.de/viewtopic,p,80448.html#80448
% Argumente: siehe http://texwelt.de/blog/modifizieren-eines-biblatex-stils/
\DeclareNameFormat{sortname}{% Bibliographie
  \ifnum\value{uniquename}=0 % Normalfall
    \ifuseprefix%
      {%
         \usebibmacro{name:family-given}
           {\namepartfamily}
           {\namepartgiveni}
           {\namepartprefix}
           {\namepartsuffixi}%
       }
      {%
         \usebibmacro{name:family-given}
           {\namepartfamily}
           {\namepartgiveni}
           {\namepartprefixi}
           {\namepartsuffixi}%
       }%
  \fi
  \ifnum\value{uniquename}=1% falls nicht eindeutig, abgek. Vorname 
      {%
         \usebibmacro{name:family-given}
           {\namepartfamily}
           {\namepartgiveni}
           {\namepartprefix}
           {\namepartsuffix}%
       }%
  \fi
  \ifnum\value{uniquename}=2% falls nicht eindeutig, ganzer Vorname 
      {%
         \usebibmacro{name:family-given}
           {\namepartfamily}
           {\namepartgiven}
           {\namepartprefix}
           {\namepartsuffix}%
       }%
  \fi   
  \usebibmacro{name:andothers}}

\DeclareNameFormat{labelname}{% für Zitate
  \ifnum\value{uniquename}=0 % Normalfall
    \ifuseprefix%
      {%
         \usebibmacro{name:family-given}
           {\namepartfamily}
           {\empty}
           {\namepartprefix}
           {\namepartsuffixi}%
       }
      {%
         \usebibmacro{name:family-given}
           {\namepartfamily}
           {\empty}
           {\namepartprefixi}
           {\namepartsuffixi}%
       }%
  \fi
  \ifnum\value{uniquename}=1% falls nicht eindeutig, abgek. Vorname 
      {%
         \usebibmacro{name:family-given}
           {\namepartfamily}
           {\namepartgiveni}
           {\namepartprefix}
           {\namepartsuffix}%
       }%
  \fi
  \ifnum\value{uniquename}=2% falls nicht eindeutig, ganzer Vorname 
      {%
         \usebibmacro{name:family-given}
           {\namepartfamily}
           {\namepartgiven}
           {\namepartprefix}
           {\namepartsuffix}%
       }%
  \fi   
  \usebibmacro{name:andothers}}
      
  
\DeclareFieldFormat{extrayear}{% = the 'a' in 'Jones 1995a'
  \iffieldnums{labelyear}
    {\mknumalph{#1}}
    {\mknumalph{#1}}}        

\renewcommand*{\multinamedelim}{\addslash}
\renewcommand*{\finalnamedelim}{\addslash}
\renewcommand*{\multilistdelim}{\addslash}
\renewcommand*{\finallistdelim}{\addslash}

\renewcommand{\nameyeardelim}{~}

% Literaturverzeichnis: Doppelpunkt zwischen Name (Jahr): Rest 
% http://de.comp.text.tex.narkive.com/Tn1HUIXB/biblatex-authoryear-und-doppelpunkt
\renewcommand{\labelnamepunct}{\addcolon\addspace}

% damit die Darstellung für Vollzitate von Primärquellen in 
% Fußnoten später auf "nicht fett" geändert werden kann 
% (nur für Zitate von Sekundärliteratur relevant)
\newcommand{\textfett}[1]{\textbf{#1}}

% für Zitate von Sekundärliteratur:
\newcommand{\footcitePrimaerSekundaer}[4]{%
  \renewcommand{\textfett}[1]{##1}%
  \footnote{\fullcite[#2]{#1}, zitiert nach \cite[#4]{#3}}%  
  \renewcommand{\textfett}[1]{\textbf{##1}}%
}

% Im Literaturverzeichnis: Autor (Jahr) fett
\renewbibmacro*{author}{%
  \ifboolexpr{%
    test \ifuseauthor%
    and
    not test {\ifnameundef{author}}
  }
    {\usebibmacro{bbx:dashcheck}
       {\bibnamedash}
       {\usebibmacro{bbx:savehash}%
        \textfett{\printnames{author}}%
        \iffieldundef{authortype}
          {\setunit{\addspace}}
          {\setunit{\addcomma\space}}}%
     \iffieldundef{authortype}
       {}
       {\usebibmacro{authorstrg}%
        \setunit{\addspace}}}%
    {\global\undef\bbx@lasthash
     \usebibmacro{labeltitle}%
     \setunit*{\addspace}}%
  \textfett{\usebibmacro{date+extrayear}}}

% Sonderfall: Quelle ohne Autor, aber mit Herausgeber
% Name des Herausgebers wird fett gedruckt
\renewbibmacro*{bbx:editor}[1]{%
  \ifboolexpr{%
    test \ifuseeditor%
    and
    not test {\ifnameundef{editor}}
  }
    {\usebibmacro{bbx:dashcheck}
       {\bibnamedash}
       {\textfett{\printnames{editor}}%
        \setunit{\addcomma\space}%
        \usebibmacro{bbx:savehash}}%
     \usebibmacro{#1}%
     \clearname{editor}%
     \setunit{\addspace}}%
    {\global\undef\bbx@lasthash
     \usebibmacro{labeltitle}%
     \setunit*{\addspace}}%
  \textfett{\usebibmacro{date+extrayear}}}

% Anpassungen für deutsche Sprache
\DefineBibliographyStrings{ngerman}{%
	nodate = {{o.J.}},
	urlseen = {{Abruf:}},
	ibidem = {{ebenda}}
}

% keine Anführungszeichen beim Titel im Literaturverzeichnis
\DeclareFieldFormat[article,book,inbook,inproceedings,manual,misc,phdthesis,thesis,online,report]{title}{#1\isdot}

\newcommand{\literaturverzeichnis}{%
% nur Literaturverzeichnis
% (als eigenes Kapitel)
\phantomsection
\addcontentsline{toc}{chapter}{Literaturverzeichnis}
\spezialkopfzeile{Literaturverzeichnis}
\defbibheading{lit}{\chapter*{Literaturverzeichnis}}
\label{chapter:quellen}
\printbibliography[heading=lit,notkeyword=ausblenden]
} % mit DHBW-spezifischen Einstellungen

\usepackage{hyperref}       % URL-Formatierung, klickbare Verweise

\usepackage{tocloft}        % für Verzeichnis der Anhänge

\newcounter{anhcnt}
\setcounter{anhcnt}{0}
\newlistof{anhang}{app}{}

\newcommand{\anhang}[1]{%
  \refstepcounter{anhcnt}
  \setcounter{anhteilcnt}{0}
  \section*{Anhang \theanhcnt: #1}
  \addcontentsline{app}{section}{\protect\numberline{Anhang \theanhcnt}#1}\par
}

\newcounter{anhteilcnt}
\setcounter{anhteilcnt}{0}

\newcommand{\anhangteil}[1]{%
	\refstepcounter{anhteilcnt}
	\subsection*{Anhang~\arabic{anhcnt}/\arabic{anhteilcnt}: #1}
	\addcontentsline{app}{subsection}{\protect\numberline{Anhang \theanhcnt/\arabic{anhteilcnt}}#1}\par
}

\renewcommand{\theanhteilcnt}{Anhang \theanhcnt/\arabic{anhteilcnt}}

% vgl. S. 4 Paket-Beschreibung tocloft 	
% Einrückungen für Anhangverzeichnis
\makeatletter
\newcommand{\abstaendeanhangverzeichnis}{
\renewcommand*{\l@section}{\@dottedtocline{1}{0em}{5.5em}}
\renewcommand*{\l@subsection}{\@dottedtocline{2}{2.3em}{6.5em}}
}
\makeatother

% Abbildungs- und Tabellenverzeichnis
% Bezeichnungen
\renewcaptionname{ngerman}{\figurename}{Abb.}
\renewcaptionname{ngerman}{\tablename}{Tab.}
% Einrückungen
\makeatletter
\renewcommand*{\l@figure}{\@dottedtocline{1}{0em}{2.3em}}
\renewcommand*{\l@table}{\@dottedtocline{1}{0em}{2.3em}}
\makeatother


\usepackage{chngcntr}                % fortlaufende Zähler für Fußnoten, Abbildungen und Tabellen
\counterwithout{figure}{chapter}
\counterwithout{table}{chapter}
\counterwithout{footnote}{chapter}

\usepackage[automark]{scrlayer-scrpage} 
%% Definitionen für Kopf- und Fußzeile auf normalen Seiten
\defpagestyle{kopfzeile}
{% Kopfdefinition
  (\textwidth,0pt)    % Länge der oberen Linie,Dicke der oberen Linie       
  {} % Definition für linke Seiten im doppelseitigen Layout
  {} % Definition für rechte Seiten im doppelseitigen Layout      
  {  % Definition für Seiten im einseitigen Layout
	\makebox[0pt][l]{\rightmark}% 
	\makebox[\linewidth]{}% 
  }        
  (\textwidth, 0.4pt) % Untere Linienlänge, Untere Liniendicke
}
{% Fußdefinition
  (\textwidth,0pt)    % Obere Linienlänge, Obere Liniendicke
  {} % Definition für linke Seiten im doppelseitigen Layout
  {} % Definition für rechte Seiten im doppelseitigen Layout
  {  % Definition für Seiten im einseitigen Layout
    \makebox[\linewidth]{}%
    \makebox[0pt][r]{\pagemark}%
  }
  (\textwidth, 0pt)   % Länge der unteren Linie,Dicke der unteren Linie
}

%% Definitionen für Kopf- und Fußzeile auf ersten Seiten eines Kapitels
\defpagestyle{kapitelkopfzeile}
{% Kopfdefinition
  (\textwidth,0pt)    % Länge der oberen Linie,Dicke der oberen Linie       
  {} % Definition für linke Seiten im doppelseitigen Layout
  {} % Definition für rechte Seiten im doppelseitigen Layout      
  {}  % Definition für Seiten im einseitigen Layout
  (\textwidth, 0pt) % Untere Linienlänge, Untere Liniendicke
}
{% Fußdefinition
  (\textwidth,0pt)    % Obere Linienlänge, Obere Liniendicke
  {} % Definition für linke Seiten im doppelseitigen Layout
  {} % Definition für rechte Seiten im doppelseitigen Layout
  {  % Definition für Seiten im einseitigen Layout
    \makebox[\linewidth]{}%
    \makebox[0pt][r]{\pagemark}%
  }
  (\textwidth, 0pt)   % Länge der unteren Linie,Dicke der unteren Linie
}

%% Definitionen für Kopf- und Fußzeile im Anhang und bei Quellenverzeichnisse
\newcommand{\spezialkopfzeileBezeichnung}{}
\defpagestyle{spezialkopfzeile}
{% Kopfdefinition
  (\textwidth,0pt)    % Länge der oberen Linie,Dicke der oberen Linie       
  {} % Definition für linke Seiten im doppelseitigen Layout
  {} % Definition für rechte Seiten im doppelseitigen Layout      
  {  % Definition für Seiten im einseitigen Layout
	\makebox[0pt][l]{\spezialkopfzeileBezeichnung}% 
	\makebox[\linewidth]{}% 
  }        
  (\textwidth, 0.4pt) % Untere Linienlänge, Untere Liniendicke
}
{% Fußdefinition
  (\textwidth,0pt)    % Obere Linienlänge, Obere Liniendicke
  {} % Definition für linke Seiten im doppelseitigen Layout
  {} % Definition für rechte Seiten im doppelseitigen Layout
  {  % Definition für Seiten im einseitigen Layout
    \makebox[\linewidth]{}%
    \makebox[0pt][r]{\pagemark}%
  }
  (\textwidth, 0pt)   % Länge der unteren Linie,Dicke der unteren Linie
}
            
\newcommand\spezialkopfzeile[1]{%
  \renewcommand\spezialkopfzeileBezeichnung{#1}
  \pagestyle{spezialkopfzeile}
}
                
% Standard-Pagestyle auswählen
\pagestyle{kopfzeile}

% keine Kopfzeile anzeigen auf Seiten, auf denen ein 
% Kapitel beginnt oder das Inhalts-/Abbildungs-/Tabellenverzeichnis steht 
\renewcommand{\chapterpagestyle}{kapitelkopfzeile}
\tocloftpagestyle{kapitelkopfzeile}

		 % für schöne Kopfzeilen 

\usepackage{textcomp}            % erlaubt EUR-Zeichen in Eingabedatei
\usepackage{eurosym}             % offizielles EUR-Symbol in Ausgabe
\renewcommand{\texteuro}{\euro}  % ACHTUNG: nach hyperref aufrufen!

\usepackage{scrhack}             % stellt Kompatibilität zw. KOMA-Script
                                 % (scrreprt) und anderen Paketen her
                                 
% Anpassung der Abstände bei Kapitelüberschriften
% (betrifft auch Inhalts-, Abbildungs- und Tabellenverzeichnis)
\renewcommand*\chapterheadstartvskip{\vspace*{-\topskip}}
\newcommand{\myBeforeTitleSkip}{1mm}
\newcommand{\myAfterTitleSkip}{10mm}
\setlength\cftbeforetoctitleskip{\myBeforeTitleSkip}
\setlength\cftbeforeloftitleskip{\myBeforeTitleSkip}
\setlength\cftbeforelottitleskip{\myBeforeTitleSkip}

\setlength\cftaftertoctitleskip{\myAfterTitleSkip}
\setlength\cftafterloftitleskip{\myAfterTitleSkip}
\setlength\cftafterlottitleskip{\myAfterTitleSkip}

% \usepackage{pgf}
\usepackage{layouts}
\usepackage{pgfplots}
\usepgfplotslibrary{external}
\tikzexternalize
%%% Ende der Präambel %%

%%% Name der eigenen Literatur-Datenbank (ggf. anpassen) %%%
\bibliography{includes/literatur-datenbank.bib}

\begin{document}
%%% Deckblatt einbinden %%% 
% Anpassungen nötig (Name, Titel etc.)
% HIER EDITIEREN: 
% Typ der Arbeit (für Deckblatt und ehrenwörtliche Erklärung)
% - bitte Zutreffendes auswählen
%\newcommand{\typMeinerArbeit}{1. Projektarbeit} 
%\newcommand{\typMeinerArbeit}{2. Projektarbeit} 
%\newcommand{\typMeinerArbeit}{Seminararbeit} 
\newcommand{\typMeinerArbeit}{Bachelorarbeit} 

% Thema der Arbeit (für ehrenwörtliche Erklärung, ohne Umbrüche)
% HIER EDITIEREN: 
\newcommand{\themaMeinerArbeit}{Existiert eine Korrelation zwischen Storypoint-Aufwandsabschätzungen und Soft-warekomplexitätsmetriken? - Eine deskriptive Fallstudie sechs agiler Softwareprojekte}

% Vorname, Name der Autorin/des Autors (für Titelseite und Metadaten)
% HIER EDITIEREN:
\newcommand{\meinName}{Tim Struthoff}

\thispagestyle{empty}

\begin{spacing}{1}
\begin{center}	
~\vspace{0mm}

% HIER EDITIEREN: Titel der Arbeit
{\sffamily
\Large
% \Large  % bei sehr langen Titeln ggf. etwas kleinere Schriftart wählen
\textbf{Existiert eine Korrelation zwischen Storypoint-Aufwandsabschätzungen und Softwarekomplexitätsmetriken? }

\bigskip
\textbf{Eine deskriptive Fallstudie sechs agiler Softwareprojekte}
}


\vspace{15mm}

% Typ wird automatisch eingefügt (oben festlegen)
{\Large \typMeinerArbeit}

\vspace{1cm}

% HIER ggf. EDITIEREN
vorgelegt am 7. Mai 2022

\vspace{15mm}

Fakultät Wirtschaft
\medskip

Studiengang Wirtschaftsinformatik
\medskip

% HIER EDITIEREN: Kurs eintragen
Kurs WI2019I 

\vspace{10mm}

von

\vspace{10mm}

% Vorname und Name wird automatisch eingefügt (oben festlegen) 
{\large\textsc{\meinName}}

\vspace{10mm}
\end{center}

\vfill

% HIER EDITIEREN: Name des Unternehmens, Name der Betreuerin/des Betreuers
\begin{tabular}{ll}
Betreuer in der Ausbildungsstätte: & DHBW Stuttgart: \\
\hspace{0.4\linewidth} & \\
DXC Technologies & Katja Sattler \\
Andreas Jordan & Delivery Lead, Testing and Digital Assurance \\
Managing Consultant \\
\\
Unterschrift der Betreuerin/des Betreuers \\
\end{tabular}


\vspace{1cm}
%(etwas Platz für die Unterschrift der Betreuerin/des Betreuers aus der Ausbildungsstätte)
\end{spacing}

% falls ein Vertraulichkeitsvermerk erforderlich ist,
% die Kommentarzeichen in den nachfolgenden Zeilen entfernen:
 
%\begin{center}
%\small
%\textbf{Vertraulichkeitsvermerk}:
%Der Inhalt dieser Arbeit darf weder als Ganzes noch in Auszügen \\
%Personen außerhalb des Prüfungs- und Evaluationsverfahrens zugänglich gemacht werden, sofern keine anders lautende Genehmigung des Dualen Partners vorliegt. 
%\end{center}

% Meta-Daten für PDF-Datei basierend auf obigen Angaben
\hypersetup{pdftitle={\themaMeinerArbeit}}
\hypersetup{pdfauthor={\meinName}}
\hypersetup{pdfsubject={\typMeinerArbeit\ DHBW Stuttgart \the\year}}

%%% Umstellung der Seiten-Nummerierung auf i, ii, iii ... %%%
\pagenumbering{Roman} 

%%% Abstract einbinden (optionale Kurzfassung Ihrer Arbeit) %%%
\begin{abstract}
\thispagestyle{kapitelkopfzeile}
\textbf{\LaTeX-Vorlage für Projekt-, Seminar- und Bachelorarbeiten}

Bei dem vorliegenden Dokument handelt es sich um eine Vorlage, die
für Projekt-, Seminar- und Bachelorarbeiten im Studiengang
Wirtschaftsinformatik der DHBW Stuttgart verwendet werden kann.

Sie setzt die technischen Vorgaben der Zitierrichtlinien\footnote{Sie finden diese unter \enquote{Prüfungsleistungen} im Studierendenportal (\url{https://studium.dhbw-stuttgart.de/winf/pruefungsleistungen/}).} des Studiengangs
(Stand: 01/2020) um.

\emph{Hinweise:} Bitte lesen Sie sich die Zitierrichtlinien unbedingt genau durch. Dieses Dokument ersetzt keine Anleitung oder Einführung in \LaTeX,
für die Nutzung sind daher gewisse Vorkenntnisse unerlässlich. Ein Einstieg in 
\LaTeX\ ist aber weniger schwierig, als es vielleicht auf den ersten Blick scheint
und lohnt sich für das Verfassen wissenschaftlicher Arbeiten in jedem Fall.\footnote{%
so auch \url{http://www.spiegel.de/netzwelt/tech/textsatz-keine-angst-vor-latex-a-549509.html}} 
Als Hilfestellung beim Schreiben eines Dokuments habe ich einen zweiseitigen kompakten \LaTeX-Spickzettel erstellt, der über Moodle verfügbar ist.

Ihre Rückmeldungen und Anregungen zu dieser Vorlage nehme ich gerne per E-Mail an die Adresse
\url{tobias.straub@dhbw-stuttgart.de} entgegen.

--- Prof. Dr. Tobias Straub

\vspace{5em}

\begin{center}\small
\begin{tabular}{ccl}
\multicolumn{3}{c}{\textbf{Versionshistorie}}\\
\hline
1.0	& 05.02.2015 & erste Fassung \\
\hline
1.1 & 16.02.2015 & siehe~\ref{anhang:ReleaseNotes11} \\
\hline
1.2 & 20.04.2015 & siehe~\ref{anhang:ReleaseNotes12} \\
\hline
1.3 & 20.02.2016 & siehe~\ref{anhang:ReleaseNotes13} \\
\hline
1.4 & 24.07.2017 & siehe~\ref{anhang:ReleaseNotes14} \\
\hline
1.5 & 07.01.2018 & siehe~\ref{anhang:ReleaseNotes15} \\
\hline
1.6 & 07.04.2018 & siehe~\ref{anhang:ReleaseNotes16} \\
\hline
1.7 & 12.02.2019 & siehe~\ref{anhang:ReleaseNotes17} \\
\hline
1.8 & 10.02.2020 & siehe~\ref{anhang:ReleaseNotes18} \\
\end{tabular}
\end{center}

\end{abstract}


\cleardoublepage

%%% Inhalts-, Abbildungs-, Tabellenverzeichnisse %%%
% sollen einzeilig gesetzt werden, um Platz zu sparen 
\begin{spacing}{1}
\tableofcontents
\clearpage
\chapter*{Abkürzungsverzeichnis}
\addcontentsline{toc}{chapter}{Abkürzungsverzeichnis}

\begin{acronym}[DHBW] 
% Argument definiert die Breite der ersten Spalte anhand des längsten vorkommenden Eintrags
\acro{CRM}{Customer Relationship Management}
\acro{DIL}{Digital Innovation Lab}
\acro{}{}
\acro{}{}
\acro{}{}
\acro{}{}


\end{acronym}


\clearpage
\thispagestyle{kapitelkopfzeile}
\listoffigures
\phantomsection
\addcontentsline{toc}{chapter}{Abbildungsverzeichnis} % Abb.verz. ins Inh.verz. aufnehmen

\clearpage
\listoftables
\phantomsection
\addcontentsline{toc}{chapter}{Tabellenverzeichnis}   % Tab.verz. ins Inh.verz. aufnehmen
\end{spacing}

%%% Umstellung der Seiten-Nummerierung auf 1, 2, 3 ... %%%
\cleardoublepage
\pagenumbering{arabic}

%%% Ihr eigentlicher Inhalt %%%
% Empfehlung: strukturieren Sie Ihren Text in einzelnen Dateien 
% und binden Sie diese hier mit \input{includes/dateiname.tex} ein

\chapter{Cheatsheet}

Text\footnote{Fußnote}.

Referenz \ref{section:zeichencodierung}

Abschnitten~\ref{section:werkzeuge}

\url{https://ctan.org/pkg/hyperref}

\verb|_dhbw_|

\begin{verbatim}
_dhbw_biblatex-config.tex  (weitere Einstellung für Biblatex)
_dhbw_erklaerung.tex       (ehrenwörtliche Erklärung)
_dhbw_kopfzeilen.tex       (Kapitelname in Kopfzeilen) 
_dhbw_praeambel.tex        (Einbindung der benötigten Pakete)
\end{verbatim}

\lstset{language=TeX} 
\begin{lstlisting}
% HIER EDITIEREN: 
\end{lstlisting}

\section{Werkzeuge}\label{section:werkzeuge}

\subsection{Titel und Erklärung}

\subsubsection{Titel und Erklärung}

\emph{Hinweis:}


\chapter{Einleitung}\label{einleitung}

Diese Bachelorarbeit wird als Teil eines Praktikums in der Digital
Innovation Lab (DIL) Abteilung des IT-Beratungs- und
Dienstleistungsunternehmens DXC Technology geschrieben. Die DXC
Technology befasst sich unter anderem mit der Herstellung von
Individualsoftware für eine Vielzahl von Kunden\footnote{Interne quelle}.
Die DIL Abteilung im Speziellen ist dabei für die Erstellung von ersten,
minimal lauffähigen Versionen (Minimum Viable Products, MVPs) dieser
Softwareprodukte verantwortlich. Diese MVPs werden dann von anderen
Abteilungen der DXC zu größeren Softwareprodukten weiterentwickelt.

Ein Teil des Dienstleistungsangebotes der DXC liegt in der Wartung,
Betreuung und Weiterentwicklung eben dieser Software. Angebote zur
Wartung und Weiterentwicklung von Software stellen auch eine wesentliche
Umsatzquelle der DXC dar. Dieser Bereich ist also aus
betriebswirtschaftlicher Sicht von besonderem Interesse. Auch aus der
Perspektive der Kunden der DXC ist der Bereich der Wartung und
Weiterentwicklung der Software von Interesse. So entfällt ein
signifikanter Anteil der Kosten eines Softwareproduktes auf deren
Wartung und Weiterentwicklung\footnote{Quelle fehlt}. Im Interesse des
Kunden sollten diese Kosten gesenkt werden\footnote{Interview Katja}.

Eine Vielzahl an Studien kamen zu dem Konsens, dass die Komplexität
eines Softwareproduktes einen wesentlichen Einflussfaktor für den
Aufwand von Wartung, Betrieb und Weiterentwicklung der Software
darstellt. Jones 2008 konnte beweisen, dass die Komplexität und der
Umfang von Software stark und direkt mit dem Wartungsaufwand\footnote{(Jones
  2008 S. 64, 335, 627)} und der durchschnittlichen Fehleranzahl
korrelieren\footnote{(Jones 2008 S. 64, 503).}. Also sei es sinnvoll,
die Komplexität der Software permanent zu beobachten\footnote{(Jones
  2008, S. 503).}.

\section{Zielsetzung}\label{zielsetzung}

Motiviert durch die betriebswirtschaftliche Relevanz der
Softwarekomplexität liegt die Zielsetzung dieser Arbeit in der Auswahl
und Validierung von Methoden zur automatisierten Bestimmung der
Komplexität von Software. Die Auswahl der Berechnungsmethoden erfolgt
anhand einer umfangreichen Literaturanalyse. Das Arbeitsergebnis ist in
diesem ersten Schritt eine Aufstellung von Berechnungsmethoden mit
jeweils einer Erklärung. So soll es den Leser*innen möglich sein, sich
ein umfangreiches Bild der aktuellen Praxis in der
Softwarekomplexitätsbestimmung zu verschaffen.

Für die Verifizierung der zuvor aufgestellten Metriken wird ein
Vergleich angestellt. Dabei werden für fünf Projekte der DXC
Technologies, sowie für ein externes Projekt Komplexitätsabschätzungen
von Experten mit den berechneten Metriken verglichen. Als Ergebnis
dieses Arbeitsschrittes ist eine Bestimmung des Grades der Korrelation
vorgesehen.

Mit diesem Arbeitsergebnis soll dann in einem letzten Schritt ein
Ausblick auf die weitere Verwendung der Ergebnisse gegeben werden.
Insbesondere wird die zukünftige Umsetzung einer sog.
Umgebungsparameteranalyse in Aussicht gestellt. Diese wurde bereits von
einem unternehmensinternen Experten skizziert. Dabei soll das Ergebnis
dieser Arbeit mit verschiedenen Einflussfaktoren der Projektumgebung
verglichen werden. Aus diesem Vergleich sollen logische Schlüsse auf
mögliche Einflussfaktoren gebildet werden.

\section{Verwandte Arbeiten}\label{verwandte-arbeiten}

Das Feld der Komplexitätsbestimmung von Software besteht bereits ähnlich
lange wie die Softwareentwicklung selbst\footnote{quelle}. Erste
Komplexitätsmaßzahlen, wie z.B. die zyklomatische Komplexität von
McCabe\footnote{McCabe, T. J. 1976} wurden in den sechziger und
siebziger Jahren des zqwanzigsten Jahrhunderts entwickelt\footnote{Zuse,
  H. 1991, S. 25} \footnote{Rubey, R. J./Hartwick, R. D. 1968}. Also ist
anzunehmen, dass auch zu der Verifizierung dieser Arbeiten bereits eine
Vielzahl an theoretischen Abhandlungen existieren. Durch eine
Literaturrecherche konnten einige Arbeiten zu der Verifizierung von
Softwarekomplexitätsmetriken identifiziert werden. Diese werden im
Folgenden zusammenfassend beschrieben.

In Kemerer\footnote{Kemerer, C. F. 1987} werden verschiedene
Komplexitätsmaßen auf 15 Projekte einer Firma angewendet. Dabei kommen
die Codezeilen basierten Maßen SLIM und COCOMO, sowie die
nicht-Codezeilen-basierten Maßen ESTIMACS und Function Points zum
Einsatz\footnote{Ebenda, S. 2}. Die Autoren kamen zu dem Ergebnis, dass
die Metriken nur aussagekräftig sind, wenn sie auf die Projekte
individuell kalibriert werden. Mit der Kalibrierung konnten
Genauigkeitsraten von 88\% erzielt werden\footnote{Ebenda, S. 12}.

In rumreichExaminingSoftwareDesign2019 werden die Komplexitätsmetriken
Anzahl Codezeilen, die zyklomatische Komplexität, die Halstead
Komplexitätsmaßzahl und der Maintainability Index auf Projekte von
Studenten angewendet, um so zu erfahren, ob eine Änderung der
Lehrmethode die Komplexität der Projekte beeinflusst\footnote{(Rumreich
  and Kecskemety 2019:1)}

In aleneziEmpiricalAnalysisComplexity2015 wird der generelle Verlauf von
Softwarekomplexitätsmetriken untersucht. Dabei kann das sechste
Lehmansche Gesetz bewiesen warden, nach dem die Komplexität einer
Software über Zeit steigt\footnote{(Alenezi and Almustafa 2015:262)}\textbf{.}

\section{Forschungsbeitrag}\label{forschungsbeitrag}

Der Forschungsbeitrag dieser Arbeit liegt in der Beschreibung der
Korrelation von mathematisch berechneten Komplexitätsmaßen mit
Aufwandsabschätzungen von Experten. Dieser Arbeitsbeitrag stellt in
dreierlei Hinsicht einen Zusatznutzen dar:

Zum einen können Softwarekomplexitätsmetriken für den spezifischen
Anwendungsbereich des DIL ausgewählt und validiert werden. Das
ermöglicht es der Abteilung, diese Maßzahlen in Zukunft für die
kontinuierliche Analyse bestehender Projekte sowie für die initiale
Beurteilung neuer Projekte zu verwenden.

Zweitens werden die Metriken nicht nur im Kontext des DIL, sondern auch
für die Allgemeinheit validiert. Zu den hier behandelten Projekten
existieren noch keine öffentlichen Softwarekomplexitätsanalysen. Mit
dieser Arbeit können die Metriken also besser beurteilt werden.

Drittens wurde im Laufe dieser Arbeit festgestellt, dass eine Abweichung
der Metriken von den Aufwandsabschätzungen ein Indikator für
Prozessereignisse der Softwareentwicklung seien könnte. Diese Arbeit
könnte also die Basis für ein System zur kontinuierlichen Verbesserung
agiler Softwareentwicklungspraktiken darstellen.

\section{Methodisches Vorgehen}\label{methodisches-vorgehen}

Als Abschlussarbeit eines Wirtschaftsinformatik (WI) Studiums wird sich
diese Arbeit auch an den Methodiken der WI orientieren. Die WI ist in
ihrer Erkenntnisgewinnung methodenpluralistisch aufgestellt\footnote{Wilde,
  T./Hess, T. o. J., S. 1}. Ihr instrumentales Portfolio beinhaltet
sowohl Methodiken aus den Real-, den Formal- sowie den
Ingenieurswissenschaften\footnote{Ebenda, S. 1}.

Unter einer Methode wird generell eine spezielle Vorgehensweise
verstanden. Sie zeichnet sich durch ein Regelsystem von
Untersuchungsinstrumenten aus. Ist diese Methode als wissenschaftlich zu
klassifizieren, müssen diese Regeln auch intersubjektiv nachvollziehbar
sein\footnote{Herrmann (1999)}.

In der WI lassen sich zwei erkenntnistheoretische Ansätze herausstellen:
Das konstruktionswissenschaftliche Paradigma strebt nach dem Schaffen
und Evaluieren von Informationssystemen in Form von Modellen, Methoden
und Softwaresystemen\footnote{(Wilde/Hess S. 2).}. Es weist z.B.
Vorgehen der Informationssystemsgestaltung, wie das Prototyping
auf\footnote{(Wilde/Hess S. 3 und Simon 1998)}. Dem gegenüber steht das
behavioristische und verhaltenswissenschaftliche Paradigma. Nach diesem
Paradigma wird in der WI das Verhalten und die Auswirkungen von
bestehenden Informationssystemen untersucht\footnote{(vgl. Wilde / Hess
  2006 S. 3)}.

In dieser Arbeit soll nach dem behavioristischen Paradigma eine
Erkenntnis gewonnen werden\footnote{(vgl. Wilde / Hess 2006 S. 3).}. Der
bestehende Sachzusammenhang der Korrelation von mathematischen
Komplexitätsmetriken mit subjektiven Aufwandsabschätzungen in agilen
Projekten soll untersucht werden. Aus dieser Untersuchung soll induktiv
auf einen Gesamtzusammenhang bzw. eine Gesetzesmäßigkeit geschlossen
werden.

Als Hypothese zu dieser Gesetzesmäßigkeit wird in dieser Arbeit eine
eingeschränkte Korrelation zwischen den Aufwandsabschätzungen und den
Codekomplexitätsmetriken vermutet. Es ist grundsätzlich anzunehmen, dass
eine Erhöhung der mathematischen Komplexität auch in einer Erhöhung des
geschätzten Aufwandes widergespiegelt wird. Jedoch sind auch
Störfaktoren dieser Korrelation abzusehen.

Die Validierung einer Metrik durch die Untersuchung ihrer Korrelation
mit einer anderen Größe ist nach Zuse, Fenton und Bowl ein verbreiteter
Ansatz\footnote{(Zuse, 1991, p. 561)} \footnote{(Zuse, 1991, p. 562)}
\footnote{Nach BOWL83 sei eine Metrik dann validiert, wenn sie das
misst, was sie angibt zu messen.}. Mit der Untersuchung dieser
Korrelation können die Komplexitätsmetriken also im Kontext des DIL
validiert werden. Eine genaue Korrelation kann und soll in dieser Arbeit
aufgrund der zu erwartenden Störfaktoren und des geringen Umfangs der
Untersuchung aber nicht berechnet werden\footnote{(Jones 2008:449)}.
Vielmehr soll eine unvollständige Theorie in Form einer Näherungsangabe
als ceteris-paribus Hypothese aufgestellt werden\footnote{(vgl. Wilde /
Hess S. 3).}.

Trotz ihres vergleichsweisen jungen Alters bietet die WI eine üppige
Bandbreite an Methoden der Erkenntnisgewinnung\footnote{heinrichForschungsmethodikIntegrationsdisziplinBeitrag2005
S. 113}. Gerade die Methodik der Fallstudienforschung erlebt einen
stetigen Anstieg in Popularität\footnote{(yinCaseStudyResearch2014 S. 22
/Michel et al., 2010 / David, 2006b, p. xxiii / David, 2006a / Mills,
Durepos, \& Wiebe, 2010a))}. Insbesondere für eine beschreibende
Forschung, wie sie in dieser Arbeit angestrebt wird, sei die
Fallstudienmethodik geeignet\footnote{(dubeRigorInformationSystems2003
S. 607).}. In einer Fallstudie wird ein Phänomen in seinem natürlichen
Kontext beschrieben. Es wird eine geringe Anzahl von Fällen intensiv
sowohl mit qualitativen als auch quantitativen Analysemethoden
untersucht\footnote{gothlichFallstudienAlsForschungsmethode2003 S. 7}.
Es werden verschiedene Datentypen gesammelt und diese in Verbindung
zueinander und zu der Hypothese gebracht. Zum Herstellen dieser
logischen Verbindungen steht eine Reihe von Werkzeugen zur Verfügung.
Insbesondere die Zeitreihenanalyse findet in dieser Arbeit Anwendung
\footnote{gothlichFallstudienAlsForschungsmethode2003 S. 6}. Der Aufbau
der Fallstudie\footnote{gothlichFallstudienAlsForschungsmethode2003 S.
8ff} wird in Kapitel 4 weiter beschrieben.

\section{Aufbau der Arbeit}\label{aufbau-der-arbeit}

Der Aufbau dieser Arbeit orientiert sich an dem Aufbau ähnlicher
Arbeiten\footnote{(Alenezi and Almustafa 2015:260), "5.1. Selected
  Systems" (Alenezi and Almustafa 2015:260), "Threats to Validity"
  (Alenezi and Almustafa 2015:264)} und soll so eine zielgerichtete und
übersichtliche Erfassung der Maßzahlen ermöglichen.

Zunächst wird das Thema der Softwarekomplexität aus einem generellen
Blickwinkel betrachtet (Kapitel 2). Dabei wird der Begriff
Softwarekomplexität zunächst definiert und dann erläutert (Kapitel 2.1
und 2.2). In einem zweiten Schritt werden verschiedene Methoden zur
Messung von Softwarekomplexität begründet ausgewählt und erklärt
(Kapitel 2.4).

Als Gegenstück zu den Softwarekomplexitätsmetriken werden in Kapitel 3
die Komplexitätsabschätzungen der Experten erläutert. Dabei wird
insbesondere auf den Kontext der Komplexitätsabschätzungen in der agilen
Entwicklung der Projekte Bezug genommen.

Nach der Definition der beiden Untersuchungsgrößen wird in Kapitel 4 der
Aufbau der Fallstudienforschung beleuchtet. Es wird ein formales
Forschungsprotokoll definiert. Insbesondere wird ein Analysealgorithmus
zur automatisierten Untersuchung der Projekte vorgestellt.

Ab Kapitel 5 findet die praktische Untersuchung der Komplexitätsmetriken
statt. Es werden die einzelnen Projekte als Fälle vorgestellt. Jeweils
wird die Datenerhebung beschrieben und mit der zuvor erläuterten
Forschungsmethodik ausgewertet. In jedem Projekt wird ein Fazit zu der
Korrelation der Metriken gezogen. In einem letzten Schritt werden
gesammelte Zusatzdaten aus den Projekten hinzugezogen, um das
Forschungsergebnis zu erklären und auch mögliche Störfaktoren
aufgezeigt.

Eine gesammelte Analyse aller Ergebnisse findet in Kapitel 6 statt. Hier
wird versucht, aus den Ergebnissen der einzelnen Projekte einen Schluss
auf eine generelle Korrelation zu ziehen.

Zuletzt wird in Kapitel 7 ein Fazit der Arbeit gegeben. Dabei wird
insbesondere auf mögliche Kritik an den Forschungsergebnissen
eingegangen (Kapitel 7.1) und ein Ausblick auf weitere Entwicklungen
gegeben (Kapitel 7.2).
\chapter{Softwarekomplexitätsmaße}\label{softwarekomplexituxe4tsmauxdfe}

Die Analyse der Softwarekomplexität ist ein Teil der statischen
Sourcecode-Analyse. In der Code-Analyse wird im Nachhinein (a
posterio\footcite[Vgl. ][S. 261]{hoffmannSoftwareQualitat2013}) anhand von statischen
Analysen einer Stichprobe zu einem bestimmten Zeitpunkt\footcite[Vgl. ][S. 86]{ebertSoftwareMetrikenPraxisEinfuhrung1996} des Sourcecodes festgestellt, ob die Software vorher
definierten Qualitätsanforderungen entspricht\footcite[Vgl. ][S. 261]{ebertSoftwareMetrikenPraxisEinfuhrung1996}.
Diese Analyse kann sowohl automatisiert als auch manuell erfolgen und
wird im Gegensatz zu Verfahren der konstruktiven Qualitätssicherung auf
bereits vorhandene Software angewendet\footcite[Vgl. ][S. 261]{hoffmannSoftwareQualitat2013}.

Als Teil der statischen Code-Analyse ist die Messung der
Softwarekomplexität ein Teil der Qualitätssicherung der
Software \footcite[Vgl. ][S. 261]{hoffmannSoftwareQualitat2013}. Dabei
sollen oft unsichtbare Eigenschaften der Software sichtbar und
quantitativ messbar gemacht werden\footcite[Vgl. ][S. 261]{hoffmannSoftwareQualitat2013}
\footcite[Vgl. ][S. 561]{zuseSoftwareComplexityMeasures1991}. Der Qualität von Software können
verschiedene Teilbereiche untergeordnet werden. Je nach Definition
gehören hierzu unter anderem Funktionalität, Laufzeit, Zuverlässigkeit,
Benutzbarkeit, Wartbarkeit, Transparenz, Übertragbarkeit und
Testbarkeit\footcite[Vgl. ][S. 22f]{hoffmannSoftwareQualitat2013} \footcite[Vgl. ][S. 245]{liggesmeyerSoftwareQualitatTestenAnalysieren2009}. Gerade auf die Qualitätsmerkmale Wartbarkeit und
Testbarkeit hat die Komplexität der Software einen direkten Einfluss. So
steigt der Aufwand des Wartens und Testens mit dem Umfang und der
Komplexität der Software.

Im Lebenszyklus einer Software kann die Komplexitätsuntersuchung als
Teil der Qualitätssicherung sowohl an das Ende der Entwicklungsphase
gestellt werden als auch kontinuierlich parallel zu der
Weiterentwicklung stattfinden.

\section{Definitionen}\label{definitionen}

Eine Betrachtung von Softwarekomplexitätsmaßen setzt zunächst eine
genaue Definition dieser voraus. Softwarekomplexitätsmetriken sind
Metriken zur Messung der Komplexität einer Software. Diese drei Teile
der Definition von Softwarekomplexitätsmetriken werden im Folgenden
definiert.

Laut dem \ac{IEEE} Standard Glossar der Softwareentwicklung besteht
\emph{Software} aus den Computerprogrammen, Prozeduren und
gegebenenfalls der Dokumentation und den Daten, die den Betrieb eines
Computersystems betreffen\footcite[Vgl. ][S. 66]{IEEEStandardGlossary}

Die Definition von \emph{Komplexität} wird allgemein als schwierig
erachtet \footcite[Vgl. ][S. 335 und 627]{jonesAppliedSoftwareMeasurement2008}.
Aber auch hier schlägt das \ac{IEEE} Standard Glossar eine allgemein
anerkannte Definition vor: „{[}Complexity is{]} the degree to which a
system or component has a design or implementation that is difficult to
understand and verify``\footcite[Vgl. ][S. 18]{IEEEStandardGlossary}. Nach dieser
Definition ist Komplexität ein Maß dafür, wie schwierig eine Software zu
verstehen und zu validieren ist. Fenton und Jones bauen auf dieser
Definition auf und schlagen vier Arten von Komplexität vor: 1. Die
Komplexität des Problems, welches die Software zu lösen versucht, 2. Die
Komplexität der Algorithmen der Software, 3. Die Komplexität der
Struktur der Software und 4. Die kognitive Komplexität der
Software \footnote{\cite[Vgl. ][S. 258]{fentonSoftwareMetricsRigorous2003} und \cite[][S. 449]{jonesAppliedSoftwareMeasurement2008}}. Nach dieser Kategorisierung bestimmen die, in dieser
Arbeit behandelten Metriken die Komplexität der Struktur der Software
(3). Weiter lässt sich zwischen inter- und intra-modularen
Komplexitätsmaßen unterscheiden. Hier werden ausschließlich
intra-modulare Komplexitätsmaßen behandelt. Diese Messen die Komplexität
einzelner Programmteile\footcite[Vgl. ][S. 7ff]{zuseSoftwareComplexityMeasures1991}.

Im Kontext der Softwaremetrie werden Metriken als ein Messystem bzw. ein
Verfahren zum Quantifizieren von Eigenschaften von Software
definiert \footnote{\cite[Vgl. ][S. 35ff]{dumkeTheorieUndPraxis1994}, \cite[][S. 4ff]{ebertSoftwareMetrikenPraxisEinfuhrung1996}, \cite[][]{augstenWasSindSoftwaremetriken} und \cite[][S. 2f]{IEEEStandardSoftware}}. Im Deutschen werden sie oft als Synonym für Maß bzw. Maßzahl
genutzt. Im Englischen findet zwischen den Begriffen „metric`` und
„measure`` eine genauere Unterscheidung statt: Eine Metrik sei eine
Funktion, die als Eingabe Daten eines Gegenstandes verwendet und hieraus
eine Zahl zur Quantifizierung dieser Eigenschaft(en) liefert\footcite[Vgl. ][S. 3]{IEEEStandardSoftware}. Ein Maß (measure) sei dahingegen die Konkrete Anwendung dieser
Metrik\footcite[Vgl. ][S. 2]{IEEEStandardSoftware}. Aufgrund dieser
Unterscheidung erscheint der Begriff Metrik für die, in dieser Arbeit
behandelten Verfahren als passender.

Zusammengefasst werden Softwarekomplexitätsmetriken in dieser Arbeit als
Verfahren zur Quantifizierung der Vielschichtigkeit der Struktur eines
Computerprogramms definiert.

\section{Kritik}\label{kritik}

Die Praktik der Softwarekomplexitätsmetrie wird allgemein stark
kritisiert.

Zum einen besteht Kritik an dem quantitativen Erfassen subjektiver
Softwareeigenschaften im Generellen. So können Metriken nicht
\emph{direkt} messen, was für Menschen als subjektive Komplexität
wahrgenommen wird, auch wenn das Verwenden von mehreren Maßzahlen hilft,
eine robustere Perspektive zu schaffen\footcite[Vgl. ][S. 2]{rumreichExaminingSoftwareDesign2019}.

Weiter werden auch die Komplexitätsmetriken im Speziellen angezweifelt:
Viele Metriken vereinen mehrere, oft konfliktäre Messziele\footcite[Vgl. ][S. 322]{fentonSoftwareMetricsRigorous2003}. Dadurch wird ihre Aussagekraft verwässert. Nicht nur an dem
mathematischen Aufbau der Metriken, sondern auch an den
Implementierungen dieser gibt es Kritik. So liefern verschiedene
Implementierungen derselben Metrik oft verschiedene
Ergebnisse\footcite[Vgl. ][S. 2]{rumreichExaminingSoftwareDesign2019}. Zusätzlich wurden
die klassischen Metriken, wie z.B. die zyklomatische Komplexität für
imperative Programmiertechniken entwickelt. Das lässt sich darauf
zurückführen, dass zu ihrem Entwicklungszeitpunkt die imperative
Programmierung noch am weitesten verbreitet war\footcite[Vgl. ][S. 277]{hoffmannSoftwareQualitat2013}. Heutzutage sind andere Programmiertechniken, wie z.B. die
objekt-orientierte Programmierung verbreiteter. Diese neuen
Programmiertechniken wurden jedoch in den klassischen Metriken nicht
berücksichtigt\footcite[Vgl. ][S. 277]{hoffmannSoftwareQualitat2013}. Ein Beispiel dieser
Problematik sind die Vererbungsmechanismen in objekt-orientierten
Programmiersprachen. Eine Vererbung erhöht z.B. in der klassischen
Metrik der zyklomatischen Komplexität nicht die Komplexität, was an
dieser Stelle die Korrelation von Code-Größe und der Komplexität des
Codes aushebelt\footcite[Vgl. ][S. 277]{hoffmannSoftwareQualitat2013}.

Ferner ist auch der Zusammenhang zwischen Softwaremetriken und externen
Softwareeigenschaften, wie z.B. der Fehleranfälligkeit
umstritten\footcite[Vgl. ][S. 627]{jonesAppliedSoftwareMeasurement2008}. Hier konnten sowohl Studien
gefunden werden, die für eine Korrelation sprechen, also auch Studien,
die dagegensprechen. In Ebert 1996 konnte z.B. ein Zusammenhang zwischen
der gemessenen Komplexität einer Aufgabe und den dabei entstandenen
Fehlern nachgewiesen werden\footcite[Vgl. ][S. 65]{ebertSoftwareMetrikenPraxisEinfuhrung1996}. Im Gegensatz dazu
konnten in Revilla 2007 keine Zusammenhänge zwischen internen
Softwaremetriken und externen Eigenschaften von Software festgestellt
werden\footcite[Vgl. ][S. 203 und 208]{revillaCorrelationsInternalSoftware2007}.

Als mögliches Fazit aus diesen Studien wird in dieser Arbeit
vorschlagen, dass die Anwendbarkeit der Maßzahlen vom jeweiligen Kontext
des Softwareprojektes und von dem erwarteten Ergebnis abhängt. Dieses
Fazit unterstreicht noch einmal die Relevanz dieser Arbeit. So kann mit
dem Ergebnis dieser Arbeit evaluiert werden, ob und welche
Softwaremaßzahlen im Kontext des \ac{DIL} Ratingen anwendbar sind.

\section{Eine Auswahl von Softwarekomplexitätsmaßzahlen}\label{eine-auswahl-von-softwarekomplexitatsmasszahlen}

Trotz der umfangreichen Kritik an der Vermessung von Software wurden in
den letzten Jahrzehnten eine Vielzahl von Maßzahlen entwickelt.

In dieser Arbeit werden insgesamt vier Metriken für die Komplexität von
Software betrachtet. Zunächst bietet die Anzahl an logischen Codezeilen
einen ersten Eindruck der Komplexität. Sie ist aber noch eher ein
Umfangsmaß als ein Komplexitätsmaß. Weiter werden die zyklomatische
Komplexität von Thomas McCabe, sowie die Softwaremaßzahlen von Halstead
behandelt. Als letzte Maßzahl wird die Einrückungskomplexität der
Autoren Hindle, Godfrey, und Holt betrachtet.

Die Auswahl der ersten drei Komplexitätsmaßen beruft sich auf ihre
generelle Popularität in der Softwaremetrie. Im Rahmen einer
umfangreichen Literaturrecherche konnten für diese Metriken die meisten
Referenzen gefunden werden. Die Maßzahl der logischen Codezeilen wird
von Zuse91 als ein wichtiges Maß zum Bestimmen des Umfangs und der
Komplexität einer Softwareanwendung eingestuft\footcite[Vgl. ][S. 145]{zuseSoftwareComplexityMeasures1991}.
Diese Position wird auch von \cite{satoExperiencesTrackingAgile2006} und
\cite{aleneziEmpiricalAnalysisComplexity2015} bekräftigt. Die zyklomatische
Komplexität von McCabe wird ebenfalls von Zuse als eine der bekanntesten
Maßzahlen eingestuft \footcite[Vgl. ][S. 145]{zuseSoftwareComplexityMeasures1991}. Weiter bekräftigt
fentonSoftwareMetricsRigorous2003, dass diese Maßzahl zum Messen der
Softwarekomplexität geeignet sei\footcite[Vgl. ][S. 31]{fentonSoftwareMetricsRigorous2003}. Auch die vier Autoren Revilla\footcite[Vgl. ][S. 203]{revillaCorrelationsInternalSoftware2007},
Jones\footcite[Vgl. ][S. 335, 627 und 449]{jonesAppliedSoftwareMeasurement2008}, Sato\footcite[Vgl. ][]{satoExperiencesTrackingAgile2006}
und Alenezi\footcite[Vgl. ][]{aleneziEmpiricalAnalysisComplexity2015} verweisen
auf die zyklomatische Komplexität als Maßzahl für die Komplexität einer
Applikation. Jones sagt dabei zusätzlich aus, dass diese
Komplexitätsmaßzahl ein besonders breites Anwendungsspektrum bedienen
könne\footcite[Vgl. ][S. 335, 627 und 449]{jonesAppliedSoftwareMeasurement2008}. Auch die Softwaremaßzahlen
von Halstead werden in der Literatur häufig referenziert. Zuse und
Revilla bekräftigen die Verbreitung dieser Maßzahlen\footnote{\cite[Vgl. ][S. 145]{zuseSoftwareComplexityMeasures1991} und \cite[S. 203]{revillaCorrelationsInternalSoftware2007}}. Zusätzlich sagt Fenton aus, dass die Maßzahlen
geeignet zum Messen der Softwarekomplexität seien\footcite[Vgl. ][S. 31]{fentonSoftwareMetricsRigorous2003}.

Zusätzlich wird die Einrückungskomplexität betrachtet, da sie einen, im
Gegensatz zu den anderen Maßzahlen methodisch sehr abweichenden Ansatz
verfolgt. So betrachtet sie im Gegensatz zu den
anderen Maßzahlen nicht die Aufteilung auf Codezeilen oder den Inhalt
des Codes, sondern die Einrückung der einzelnen Codezeilen.

Neben diesen vier Maßzahlen bestehen auch noch eine Vielzahl an weiteren
Messungsmethoden, die in dieser Arbeit aber keine weitere Betrachtung
finden sollen. Zum Beispiel wurden die gewichtete Anzahl an Methoden pro
Klasse (WMC)\footcite[Vgl. ][]{satoExperiencesTrackingAgile2006}, die Anzahl an
Kommentaren\footcite[Vgl. ][S. 207]{revillaCorrelationsInternalSoftware2007}, die Größe der Klasse\footcite[Vgl. ][]{satoExperiencesTrackingAgile2006}
und die Anzahl an Testcodezeilen\footcite[Vgl. ][]{satoExperiencesTrackingAgile2006}
nicht weiter betrachtet.

Zwischen drei der ausgewählten Maßzahlen wurden bereits Korrelationen
nachgewiesen. Zunächst besteht eine starke Korrelation zwischen der
Anzahl an logischen Codezeilen und dem Aufwand nach Halstead\footcite[Vgl. ][S. 627]{jonesAppliedSoftwareMeasurement2008}. Eine noch stärkere, linear-stabile Korrelation lässt sich
zwischen der Anzahl logischer Code Zeilen und der zyklomatischen
Komplexitätsmaßzahl nachweisen. Diese Korrelation konnte über
Programmierer*innen, Programme, Sprachen und Programmierparadigma hinweg
bewiesen werden\footnote{\cite[Vgl. ][S. 627]{jonesAppliedSoftwareMeasurement2008} und \cite[S. 137]{jayCyclomaticComplexityLines2009}}.

Die vier Komplexitätsmetriken werden nun der Reihe nach genauer
erläutert.

\subsection{Logische Codezeilen}\label{logische-codezeilen}

Als erste hier vorgestellte Metrik weist die \ac{SLOC} den geringsten Berechnungsaufwand auf. Sie wird
teilweise auch als \ac{LOC} bezeichnet\footcite[Vgl. ][S. 263]{hoffmannSoftwareQualitat2013}. Dabei werden die Zeilen an Sourcecode in der Software
berechnet\footcite[Vgl. ][S. 2]{rumreichExaminingSoftwareDesign2019}. Die Metrik lässt
sich somit ohne die Unterstützung komplexer Algorithmen berechnen und
lässt sich auf nahezu alle Programmiersprachen anwenden\footcite[Vgl. ][S. 263]{hoffmannSoftwareQualitat2013}. Eine Ausnahme bilden hier grafische
Programmiersprachen.

In ihrer Aussagekraft wird die Anzahl von Codezeilen allgemein als
begrenzt eingestuft\footnote{\cite[Vgl. ][S. 264]{hoffmannSoftwareQualitat2013} 
und \cite[S. 2]{rumreichExaminingSoftwareDesign2019}}. Unter anderem der Programmierstil und die Programmiersprache
haben einen Einfluss auf die Anzahl der Codezeilen (Hoffmann 2013:264,
Hoffmann 2013:264). Gleichzeitig erfüllt sie jedoch die Assoziativität,
Kommutativität und Monotonität\footcite[Vgl. ][S. 142]{zuseSoftwareComplexityMeasures1991} als wesentliche Qualitätskriterien für Maßzahlen.

Eine Weiterentwicklung der LOC-Metrik ist u.a. die NCSS-Metrik (Non
Commented Source Statements). Sie misst genauso wie die LOC-Metrik die
Anzahl an Codezeilen, ignoriert dabei aber alle Kommentarzeilen\footcite[Vgl. ][S. 264]{hoffmannSoftwareQualitat2013}. Zusätzlich bestehen auch empirisch ermittelte
Sprachfaktoren, die es ermöglichen, die Ergebnisse von LOC und NCSS
Metriken vor ihrer Weiterverarbeitung zu gewichten\footcite[Vgl. ][S. 264]{hoffmannSoftwareQualitat2013}.

\subsection{Zyklomatische Komplexität}\label{Zyklomatische-Komplexitat}

Die zyklomatische Komplexität wurde von Tom McCabe entwickelt und zuerst
in seinem Aufsatz „A Complexity Metric``\footcite[Vgl. ][]{mccabeComplexityMeasure1976}
veröffentlicht\footcite[Vgl. ][S. 185]{sneedSoftwareZahlenVermessung2010}. Die zyklomatische
Komplexität gehört zu den Kontrollflussmetriken\footcite[Vgl. ][S. 88]{ebertSoftwareMetrikenPraxisEinfuhrung1996}. In der Praxis wird sie von vielen Firmen
entwicklungsbegleitend kontinuierlich erhoben, um vor problematischen
Programmkomponenten frühzeitig zu warnen. Dabei fließt sie oft auch in
die Abnahmekriterien für Software ein\footcite[Vgl. ][S. 275]{hoffmannSoftwareQualitat2013}.

Die zyklomatische Komplexität ist definiert als eine Maßzahl für die
Kontrollflusskomplexität eines Programmes\footnote{\cite[Vgl. ][S. 2]{rumreichExaminingSoftwareDesign2019} und \cite[S. 335]{jonesAppliedSoftwareMeasurement2008}}. Sie misst die Anzahl von linear
unabhängigen Pfaden auf dem Kontrollflussgraphen eines Programmes und
liefert so ein normalisiertes Komplexitätsmaß\footcite[Vgl. ][S. 2]{rumreichExaminingSoftwareDesign2019}.

Zur Berechnung der zyklomatischen Komplexität wird der Programmcode
zunächst in einen Graphen aus Knoten und Kanten umgewandelt\footnote{\cite[Vgl. ][S. 185]{sneedSoftwareZahlenVermessung2010} und \cite[S. 335]{jonesAppliedSoftwareMeasurement2008}}. Dazu wird der Quelltext in einen
Syntaxbaum umgewandelt. Diese Baumstruktur spiegelt die Struktur des
Codes wider. Aus dem Syntaxbaum wird ein Kontrollflussgraph der
Anwendung konstruiert. Der Kontrollflussgraph besteht aus den
Anweisungen und Verzweigungen des Programms und stellt alle möglichen
Ablaufwege durch das Programm dar\footcite[Vgl. ][]{rackeKontrollflussdiagramme2017}.
Der Umwandlungsprozess wird in Abbildung .. beschrieben.

GRAFIK

Aus den Eigenschaften dieses Kontrollflussgraphen kann nun die
zyklomatische Komplexität mit Mitteln der Graphentheorie\footcite[Vgl. ][S. 273]{hoffmannSoftwareQualitat2013} berechnet werden. Als relevante Größen kommen dabei die
Anzahl der Kanten des Graphen (E), die Anzahl der Knoten (N), sowie die
Anzahl unabhängiger Teilgraphen (p) zum Einsatz. Die Anzahl unabhängiger
Teilgraphen entspricht der Anzahl unabhängiger Module bzw. Funktionen in
dem Programm. Wird ein Programm mit mehreren Modulen untersucht, werden
die Kanten und Knoten der Kontrollflussgraphen der Module
aufsummiert\footcite[Vgl. ][S. 275f]{hoffmannSoftwareQualitat2013}.

Die Formel für die zyklomatische Komplexität lautet wie folgt:

V(Z) = \textbar E\textbar-\textbar N\textbar{} + 2p

Die zyklomatische Komplexität (V(Z)) entspricht also der Anzahl an
Kanten (E) minus der Anzahl an Knoten (N) plus der doppelten Anzahl der
zusammenhängenden Einzelgraphen (p) (Hoffmann 2013:273, Jones 2008:335,
Hoffmann 2013:273, Ebert 1996:88, Sneed et al 2010:185). Also ist sie
stets um eins größer als die Anzahl an Verzweigungen in einem Programm
(Hoffmann 2013:274).

In der konkreten Anwendung wird die zyklomatische Komplexität oft als
Indikator für die Wartbarkeit des Code und für die Produktivität in der
Weiterentwicklung des Codes verwendet\footcite[Vgl. ][S. 336]{jonesAppliedSoftwareMeasurement2008}.
Zusätzlich liefert sie die obere Grenze der Anzahl an Tests, die für
eine vollständige Testabdeckung des Quelltextes benötigt werden.

Die zyklomatische Komplexität wurde in der Fachliteratur vielfach
zitiert und vor allem auch vielfach kritisiert.

Zunächst wird allgemein kritisiert, dass die Maßzahl nur die Komplexität
der Ablauflogik des Codes und nicht die Komplexität des gesamten Codes
misst\footcite[Vgl. ][S. 186]{sneedSoftwareZahlenVermessung2010}. Interne Eigenschaften der Knoten\footnote{\cite[Vgl. ][S. 73]{hoffmannSoftwareQualitat2013} und \cite[S. 2]{rumreichExaminingSoftwareDesign2019}}, die Komplexität im Datenfluss\footcite[Vgl. ][S. 2]{rumreichExaminingSoftwareDesign2019} sowie die Verschachtelung von Modulen werden
ignoriert\footcite[Vgl. ][S. 89ff]{zuseSoftwareComplexityMeasures1991}. Also lässt sich sagen, dass die
Maßzahl für prozeduralen Code eine gewisse Bedeutung hat, in
objektorientiertem Code jedoch nur die Komplexität der einzelnen
Methodenkörper misst\footcite[Vgl. ][S. 186]{sneedSoftwareZahlenVermessung2010}.

Insgesamt wird geschlussfolgert, dass die zyklomatische Komplexität nur
mit Vorsicht anzuwenden sei und in Relation zu anderen Maßzahlen gesetzt
werden müsse\footcite[Vgl. ][S. 52]{fentonSoftwareMetricsRigorous2003}. Genau
aus diesem Grund wird sie in dieser Arbeit in den Kontext der agilen
Aufwandsabschätzungen und in Relation zu anderen Maßzahlen gesetzt.

\subsection{Halstead Metriken}\label{Halstead-Metriken}

Als dritte Komplexitätsmetrik werden die Halstead Metriken herangezogen.
Bei diesen Metriken wendet der Autor Maurice Howard Halstead\footcite[Vgl. ][]{halsteadElementsSoftwareScience1979}
Elemente der Kommunikationstheorie auf die Programmierung an\footcite[Vgl. ][S. 183]{sneedSoftwareZahlenVermessung2010}.

Die Berechnung der Halstead Metriken wird im folgenden an einem Beispiel
erklärt\footcite[Vgl. ][S. 184]{sneedSoftwareZahlenVermessung2010}. Das Beispiel behandelt
folgenden Code:

\lstset{language=C}
\begin{lstlisting}
main()
{
int countOne, countTwo, countThree, average;
scanf(``%d %d %d'', &countOne, &countTwo, &countThree);
average = (countOne + countTwo + countThree) / 3;
printf(``average = \%d'', average);
}
\end{lstlisting}


Für die Berechnung seiner Metriken zählt Halstead zunächst die
Grundelemente der Programmiersprache, nämlich die Operatoren und
Operanden\footnote{\cite[Vgl. ][S. 183]{sneedSoftwareZahlenVermessung2010} und \cite[S. 2]{rumreichExaminingSoftwareDesign2019}}. Operatoren führen Operationen auf Operanden durch. Beispiele
für Operatoren sind arithmetische Operatoren (z.B. Addition)
Vergleichsoperatoren und Zuweisungsoperatoren. Operanden sind in der
Regel Datenelemente wie Zahlen oder Text.

In diesem Beispiel sind die unterschiedlichen Operatoren: main, (),
\{\}, int, scanf, \&, =, +, /, printf, ,, ;

Die unterschiedlichen Operanden sind: countOne, countTwo, countThree,
average, "\%d \%d \%d", 3, "avg = \%d"

Zu den Operatoren und Operanden werden folgende Zahlen erfasst:

\begin{itemize}
\item
  n1: Die Anzahl an unterschiedlichen Operatoren (12)
\item
  n2: Die Anzahl an unterschiedlichen Operanden (7)
\item
  N1: Die Anzahl insgesamt vorkommender Operatoren (27)
\item
  N2: Die Anzahl insgesamt vorkommender Operanden (15)
\end{itemize}

Aus diesen Zahlen können nun die Halstead Metriken berechnet werden:

Der Wortschatz eines Programms ergibt sich aus der Summe der
unterschiedlichen Operatoren und Operanden:

\emph{Wortschatz (n) = Operatoren (n1) + Operanden (n2) 19 = 12 + 7}

Die Länge des Programmes ergibt sich aus der Summe der insgesamt
vorkommenden Operanden und Operatoren :

\emph{Länge (N) = Operatorenvorkommnisse (N1) + Operandenvorkommnisse
(N2) 42 = 27 + 15}

Weiter Berechnet Halstead die Größe des Programmcodes durch einen
Logarithmus:

\emph{Größe (Volume / V) = Länge (N) * log2 (Wortschatz (n) ) 53,71 = 42
* log(19)}

Für die Sprachkomplexität setzt er jeweils die Operatoren ins Verhältnis
zu den Operatorenvorkommnissen und die Operanden ins Verhältnis zu den
Operandenvorkommnissen. Beide Werte werden miteinander multipliziert:

\emph{Komplexität = (Operatoren / Operatorenvorkommnisse) * (Operanden /
Operandenvorkommnisse) 0,207 =ca (12 / 27) * (7 / 15)}

Zuletzt berechnet er den Aufwand als Quotient aus Programmgröße und
Komplexität. Je komplexer oder größer ein Programm ist, desto länger
dauere, es dieses Programm zu schreiben:

\emph{Aufwand (E) = Größe / Komplexität 259,47 = 53,71 / 0,207}

Die von Halstead vorgeschlagenen Metriken wurden vielfach kritisiert. Er
hat die Berechnungen nie empirisch untermauert und in späteren
empirischen Studien wurden sie sogar widerlegt\footcite[Vgl. ][S. 185]{sneedSoftwareZahlenVermessung2010}. Insbesondere sei zu kritisieren, dass die Codegröße,
Komplexität und der Aufwand nach Halstead nicht das Kriterium der
Monotonität erfüllen. Sie können also nicht auf einer Verhältnisskala
verwendetet werden und nur die Berechnung von Medianwerten sei
sinnvoll\footcite[Vgl. ][S. 142]{zuseSoftwareComplexityMeasures1991}. Eine zusätzliche Einschränkung
sei, dass die Halstead Metriken die Berechnungskomplexität in den
Vordergrund rücken und keine Aussage über den Kontrollfluss der
Anwendung treffen\footcite[Vgl. ][S. 2]{rumreichExaminingSoftwareDesign2019}. Insgesamt seien sie also kritisch anzusehen\footcite[Vgl. ][S. 185]{sneedSoftwareZahlenVermessung2010}.


\subsection{Einrückungskomplexität}\label{Einruckungskomplexitat}

Als letztes Komplexitätsmaß wird hier die Einrückungskomplexität
vorgestellt. Sie wurde von Abram Hindle, Michael W. Godfrey und Richard
C. Holt an der University of Waterloo entwickelt und zuerst mit ihrer
Arbeit „Reading Beside the Lines: Indentation as a Proxy for Complexity
Metrics `` auf der sechzehnten internationalen IEEE Konferenz zum
Verständnis von Computerprogrammen vorgestellt\footcite[Vgl. ][S. 1]{hindleReadingLinesIndentation2008}. Die Komplexitätsmaßzahl ist im Vergleich zu den anderen hier
behandelten Metriken verhältnismäßig neu, wurde aber bereits in Arbeiten
von Lalouche et. al\footcite[Vgl. ][S. 10]{gilWhenSoftwareComplexity2016}
\footcite[Vgl. ][S. 7]{gilCorrelationSizeMetric2017} und Landman et. al
\footcite[Vgl. ][S. 6]{landmanEmpiricalAnalysisRelationship2016} referenziert.

Die Einrückungskomplexität soll nun erklärt werden. In den meisten
Programmiersprachen werden bestimmte Zeilen zur besseren Lesbarkeit des
Codes eingerückt. Die Einrückungskomplexität macht sich genau diesen
Umstand zu Nutzen und verwendet die Einrückungen der Codezeilen als
Indikator für Komplexität\footcite[Vgl. ][S. 1]{hindleReadingLinesUsing2009}.

In prozeduralen Code, wie z.B. in C, kann die Einrückung von Codezeilen
Kontrollstrukturen, wie Verzweigungen und Schleifen anzeigen. In
objektorientierten Sprachen, wie C++, Java und JavaScript kann die
Einrückung eine Verkapselung in Form von Klassen, Subklassen oder
Methoden anzeigen. In funktionalen Sprachen wie OCaml und Lisp zeigt die
Einrückung einen neuen Handlungskontext, neue Funktionen oder einen
neuen Ausdruck. In jedem dieser drei Typen von
Programmiersprachen sind in der Regel Stellen im Code eingerückt, die
die Komplexität des Codes erhöhen. Also lasse sich schlussfolgern, dass
die Menge der Einrückungen ein Maß der Komplexität eines Programmes sei.

Zur Verifizierung dieser These setzten die Autoren die
Einrückungskomplexität in ein Verhältnis zu den älteren Metriken wie der
zyklomatischen Komplexität und den Halstead Metriken. Bei der
zyklomatischen Komplexität erhöhen z.B. Verzweigungen die Komplexität.
Diese werden in den meisten Programmiersprachen durch eine Einrückung
der Codezeilen gekennzeichnet. Also stünde die Vermutung einer
Korrelation der Metriken nahe\footcite[Vgl. ][S. 2]{hindleReadingLinesUsing2009}.
Diese Korrelation konnte von den Autoren in einer Studie von 278
populären Open Source Projekten erfolgreich nachgewiesen
werden. Damit sei die Validität der
Einrückungskomplexität als Komplexitätsmaß von Software nachgewiesen.

Für die Berechnung der Einrückungskomplexität werden die Einrückungen
der einzelnen Codezeilen im Quelltext betrachtet. Die Anzahl an
physischen Einrückungen in einer Zeile bezeichnet zunächst die Anzahl an
Leerzeichen bzw. Tabs, welche sich am Anfang einer Zeile befinden. Diese
werden für jede Zeile berechnet. Anhand dieser Angaben wird auf der
Ebene einer Datei berechnet, wie viele Leerzeichen bzw. Tabs einer
logischen Einrückungsebene entsprechen\footcite[Vgl. ][S. 5]{hindleReadingLinesUsing2009}. In den meisten Projekten entsprechen vier Leerzeichen bzw. ein
Tab einem Einrückungslevel\footcite[Vgl. ][S. 9]{hindleReadingLinesUsing2009}.
Dieser Zusammenhang wird von den Autoren als das Einrückungsmodell einer
Datei bezeichnet. Mit dem Einrückungsmodell kann nun die Anzahl an
logischen Einrückungen für jede Datei berechnet werden. Die Anzahl
dieser Einrückungen gibt für jede Zeile die Einrückungsebene an. Die
Summe der logischen Einrückungen wird als die Einrückungskomplexität
bezeichnet.

Ein wesentlicher Vorteil der Einrückungskomplexität ist ihre
Sprachenunabhängigkeit. Im Gegensatz zu der zyklomatischen Komplexität
und den Halstead Metriken setzt sie kein Verständnis der Grammatik einer
Programmiersprache voraus\footcite[Vgl. ][S. 1]{hindleReadingLinesUsing2009}
\footcite[Vgl. ][S. 2]{hindleReadingLinesUsing2009}.

Als zusätzlichen Vorteil führen die Autoren auf, dass die
Einrückungskomplexität weniger Berechnungsschritte benötige als andere
Komplexitätsmaßen. Demnach sei sie günstiger in der Durchführung als
manche andere Metriken\footcite[Vgl. ][S. 20f]{hindleReadingLinesUsing2009}.
Vergleichsoperatoren und Zuweisungsoperatoren. Operanden sind in der
Regel Datenelemente wie Zahlen oder Text.

In diesem Beispiel sind die unterschiedlichen Operatoren: main, (),
\{\}, int, scanf, \&, =, +, /, printf, ,, ;

Die unterschiedlichen Operanden sind: countOne, countTwo, countThree,
average, "\%d \%d \%d", 3, "avg = \%d"

Zu den Operatoren und Operanden werden folgende Zahlen erfasst:

\begin{itemize}
\item
  n1: Die Anzahl an unterschiedlichen Operatoren (12)
\item
  n2: Die Anzahl an unterschiedlichen Operanden (7)
\item
  N1: Die Anzahl insgesamt vorkommender Operatoren (27)
\item
  N2: Die Anzahl insgesamt vorkommender Operanden (15)
\end{itemize}

Aus diesen Zahlen können nun die Halstead Metriken berechnet werden:

Der Wortschatz eines Programms ergibt sich aus der Summe der
unterschiedlichen Operatoren und Operanden:

\emph{Wortschatz (n) = Operatoren (n1) + Operanden (n2) 19 = 12 + 7}

Die Länge des Programmes ergibt sich aus der Summe der insgesamt
vorkommenden Operanden und Operatoren :

\emph{Länge (N) = Operatorenvorkommnisse (N1) + Operandenvorkommnisse
(N2) 42 = 27 + 15}

Weiter Berechnet Halstead die Größe des Programmcodes durch einen
Logarithmus:

\emph{Größe (Volume / V) = Länge (N) * log2 (Wortschatz (n) ) 53,71 = 42
* log(19)}

Für die Sprachkomplexität setzt er jeweils die Operatoren ins Verhältnis
zu den Operatorenvorkommnissen und die Operanden ins Verhältnis zu den
Operandenvorkommnissen. Beide Werte werden miteinander multipliziert:

\emph{Komplexität = (Operatoren / Operatorenvorkommnisse) * (Operanden /
Operandenvorkommnisse) 0,207 =ca (12 / 27) * (7 / 15)}

Zuletzt berechnet er den Aufwand als Quotient aus Programmgröße und
Komplexität. Je komplexer oder größer ein Programm ist, desto länger
dauere, es dieses Programm zu schreiben:

\emph{Aufwand (E) = Größe / Komplexität 259,47 = 53,71 / 0,207}

Die von Halstead vorgeschlagenen Metriken wurden vielfach kritisiert. Er
hat die Berechnungen nie empirisch untermauert und in späteren
empirischen Studien wurden sie sogar widerlegt\footcite[Vgl. ][S. 185]{sneedSoftwareZahlenVermessung2010}. Insbesondere sei zu kritisieren, dass die Codegröße,
Komplexität und der Aufwand nach Halstead nicht das Kriterium der
Monotonität erfüllen. Sie können also nicht auf einer Verhältnisskala
verwendetet werden und nur die Berechnung von Medianwerten sei
sinnvoll\footcite[Vgl. ][S. 142]{zuseSoftwareComplexityMeasures1991}. Eine zusätzliche Einschränkung
sei, dass die Halstead Metriken die Berechnungskomplexität in den
Vordergrund rücken und keine Aussage über den Kontrollfluss der
Anwendung treffen\footcite[Vgl. ][S. 2]{rumreichExaminingSoftwareDesign2019}. Insgesamt seien sie also kritisch anzusehen\footcite[Vgl. ][S. 185]{sneedSoftwareZahlenVermessung2010}.


\subsection{Einrückungskomplexität}\label{Einruckungskomplexitat}

Als letztes Komplexitätsmaß wird hier die Einrückungskomplexität
vorgestellt. Sie wurde von Abram Hindle, Michael W. Godfrey und Richard
C. Holt an der University of Waterloo entwickelt und zuerst mit ihrer
Arbeit „Reading Beside the Lines: Indentation as a Proxy for Complexity
Metrics `` auf der sechzehnten internationalen IEEE Konferenz zum
Verständnis von Computerprogrammen vorgestellt\footcite[Vgl. ][S. 1]{hindleReadingLinesIndentation2008}. Die Komplexitätsmaßzahl ist im Vergleich zu den anderen hier
behandelten Metriken verhältnismäßig neu, wurde aber bereits in Arbeiten
von Lalouche et. al\footcite[Vgl. ][S. 10]{gilWhenSoftwareComplexity2016}
\footcite[Vgl. ][S. 7]{gilCorrelationSizeMetric2017} und Landman et. al
\footcite[Vgl. ][S. 6]{landmanEmpiricalAnalysisRelationship2016} referenziert.

Die Einrückungskomplexität soll nun erklärt werden. In den meisten
Programmiersprachen werden bestimmte Zeilen zur besseren Lesbarkeit des
Codes eingerückt. Die Einrückungskomplexität macht sich genau diesen
Umstand zu Nutzen und verwendet die Einrückungen der Codezeilen als
Indikator für Komplexität\footcite[Vgl. ][S. 1]{hindleReadingLinesUsing2009}.

In prozeduralen Code, wie z.B. in C, kann die Einrückung von Codezeilen
Kontrollstrukturen, wie Verzweigungen und Schleifen anzeigen. In
objektorientierten Sprachen, wie C++, Java und JavaScript kann die
Einrückung eine Verkapselung in Form von Klassen, Subklassen oder
Methoden anzeigen. In funktionalen Sprachen wie OCaml und Lisp zeigt die
Einrückung einen neuen Handlungskontext, neue Funktionen oder einen
neuen Ausdruck. In jedem dieser drei Typen von
Programmiersprachen sind in der Regel Stellen im Code eingerückt, die
die Komplexität des Codes erhöhen. Also lasse sich schlussfolgern, dass
die Menge der Einrückungen ein Maß der Komplexität eines Programmes sei.

Zur Verifizierung dieser These setzten die Autoren die
Einrückungskomplexität in ein Verhältnis zu den älteren Metriken wie der
zyklomatischen Komplexität und den Halstead Metriken. Bei der
zyklomatischen Komplexität erhöhen z.B. Verzweigungen die Komplexität.
Diese werden in den meisten Programmiersprachen durch eine Einrückung
der Codezeilen gekennzeichnet. Also stünde die Vermutung einer
Korrelation der Metriken nahe\footcite[Vgl. ][S. 2]{hindleReadingLinesUsing2009}.
Diese Korrelation konnte von den Autoren in einer Studie von 278
populären Open Source Projekten erfolgreich nachgewiesen
werden\footcite[][]{hindleReadingLinesUsing2009}. Damit sei die Validität der
Einrückungskomplexität als Komplexitätsmaß von Software nachgewiesen.

Für die Berechnung der Einrückungskomplexität werden die Einrückungen
der einzelnen Codezeilen im Quelltext betrachtet. Die Anzahl an
physischen Einrückungen in einer Zeile bezeichnet zunächst die Anzahl an
Leerzeichen bzw. Tabs, welche sich am Anfang einer Zeile befinden. Diese
werden für jede Zeile berechnet. Anhand dieser Angaben wird auf der
Ebene einer Datei berechnet, wie viele Leerzeichen bzw. Tabs einer
logischen Einrückungsebene entsprechen\footcite[Vgl. ][S. 5]{hindleReadingLinesUsing2009}. In den meisten Projekten entsprechen vier Leerzeichen bzw. ein
Tab einem Einrückungslevel\footcite[Vgl. ][S. 9]{hindleReadingLinesUsing2009}.
Dieser Zusammenhang wird von den Autoren als das Einrückungsmodell einer
Datei bezeichnet. Mit dem Einrückungsmodell kann nun die Anzahl an
logischen Einrückungen für jede Datei berechnet werden. Die Anzahl
dieser Einrückungen gibt für jede Zeile die Einrückungsebene an. Die
Summe der logischen Einrückungen wird als die Einrückungskomplexität
bezeichnet.

Ein wesentlicher Vorteil der Einrückungskomplexität ist ihre
Sprachenunabhängigkeit. Im Gegensatz zu der zyklomatischen Komplexität
und den Halstead Metriken setzt sie kein Verständnis der Grammatik einer
Programmiersprache voraus \footnote{\cite[Vgl.][S. 1]{hindleReadingLinesUsing2009} und \cite[S. 2]{hindleReadingLinesUsing2009}}.

Als zusätzlichen Vorteil führen die Autoren auf, dass die
Einrückungskomplexität weniger Berechnungsschritte benötige als andere
Komplexitätsmaßen. Demnach sei sie günstiger in der Durchführung als
manche andere Metriken\footcite[Vgl. ][S. 20f]{hindleReadingLinesUsing2009}.
\blinddocument
\chapter{Zitieren}\label{chapter:zitate}

Der Zitierstil ist so angepasst, dass er den Zitierrichtlinien des Studiengangs Wirtschaftsinformatik der DHBW Stuttgart entspricht. 

\section{Zitate in den Text einfügen}
In \LaTeX\ wird mit den Befehlen \verb|\footcite| 
oder 
\verb|\cite|
eine Referenz im Text eingefügt. Meist wird \verb|\cite| nur \emph{innerhalb} einer Fußnote benutzt. 
Damit ein vorangestelltes \enquote{Vgl.} in der Fußnote erscheint, können Sie wie folgt zitieren:
\begin{verbatim}
\footcite[Vgl.][S. 3]{Autor}
\footcite[Vgl.][]{Autor}
\end{verbatim}

Das erste optionale Argument von \verb|\footcite| wird dem Zitat vorangestellt, das zweite ist die Seitenzahl. Den selben Effekt hätte
\begin{verbatim}
\footnote{Vgl. \cite[S. 3]{Autor}}
\footnote{Vgl. \cite{Autor}}
\end{verbatim}

Hinweis: Falls \enquote{Vgl.}, aber keine Seitenzahl angeben werden soll, muss das zweite Argument vorhanden (jedoch leer) sein, ansonsten wird \enquote{Vgl.} als Seitenzahl interpretiert. Falsch ist also: 
\begin{verbatim}
\footcite[Vgl.]{Autor}  % so nicht!
\end{verbatim}



\subsection{Beispiele}
Nachfolgend ein paar Beispiele, um die korrekte Darstellung zu überprüfen:

\begin{itemize}
\item \cite{Schlosser} ist ein Buch über \LaTeX.
\item Zur Vorlesung \emph{Logik und Algebra} gibt es das gleichnamige Lehrbuch.\footcite{Staab}
\item nochmal dasselbe Buch\footcite{Staab}
\item ein weiteres Buch desselben Autors\footcite{BuschlingerStaab}

\item Der Konferenzbeitrag \cite{Ancuti} beschäftigt sich mit Bildverarbeitung.
\item Cloud Computing wird in einer Diplomarbeit erklärt.\footcite[S.~14]{Boettger:Diplomarbeit}
\item Preiß\footcite{Preiss} gibt eine Einführung in Datenbanken.
\item Eine Erläuterung, was \enquote{Intangibles} sind, findet sich bei Stoi\footcite[S.~82]{Stoi}.
\item weitere Ausführung in derselben Quelle\footcite[Vgl.][S.~84]{Stoi}
\item Laut Wikipedia\footcite[Abschnitt~5]{wiki:Wirtschaftsinformatik} ist Wirtschaftsinformatik ein interessantes Studienfach.
\item ITIL-Prozesse kann man tatsächlich auch mit \LaTeX\ dokumentieren.\footcite{Carvalho:PJ:2012-1}
\item Open-Source und Cloud-Computing in einem Buchbeitrag\footcite{Wind}
\item Buch mit zwei Autoren\footcite{MitZweiAutoren}
\item Buch mit drei Autoren\footcite{MitDreiAutoren}
\item Buch ohne Autor\footcite{OhneAutoren}
\item Buch ohne Autor und ohne Jahr\footcite{OhneAutorenOhneJahr}
\item und noch ein anderes Buch ohne Autor und ohne Jahr\footcite{OhneAutorenOhneJahr2}
\item Buch ohne Autor, aber dafür mit Herausgeber\footcite{keinAutorAberHerausgeber}

\item manche Bachelorarbeit baut auf einer vorhergehenden Projektarbeit\footcite{mayer:PA1} auf
\item das Handbuch zu BibLaTeX\footcite{biblatex:manual} und eines zu Windows 8\footcite{Win8}
\item zwei Beiträge zu Büchern\footcite{Trautwein:Nokia}${}^{,}$\footcite{Mann} und zu einem Konferenzband\footcite{Trautwein:Erfolgsfaktoren}
\item eine Online-Quelle\footcite{SAP:HANA}
\item eine plagiierte Dissertation,\footcite{GuttenPlag} nicht zur Nachahmung empfohlen
\item zum Testen, ob Umlaute und Sonderzeichen korrekt wiedergegeben werden\footcite{Umlauttest}
\end{itemize}

\subsection{Spezialfälle}
\begin{itemize}
\item \emph{Zwei Quellen am Satzende} werden durch Komma getrennt.\footcite{Staab}${}^{,}$\footcite{mayerLukas:PA1} Hier muss \verb|${}^{,}$| eingeschoben werden.
\item \emph{Eindeutigkeit:} Normalerweise wird kein Vorname des Autors angegeben. Falls es allerdings zur Eindeutigkeit\footnote{\cite{trautwein2011unternehmensplanspiele} vs. \cite{hitzler2011optimierung}} (bei gleicher Jahreszahl) erforderlich ist, wird der Vorname abgekürzt bzw.\ nötigenfalls sogar ganz ausgeschrieben mit angegeben.\footnote{Vgl. \cite{mayer:PA1} und \cite{mayerLukas:PA1}}
 
Welch ein Glück, dass Sie sich darum dank \LaTeX\ gar nicht kümmern müssen (arme Word\texttrademark-User ;-).

\item Die Verwendung von \emph{Sekundärliteratur}\footcitePrimaerSekundaer{Primaerquelle}{}{Sekundaerquelle}{11} wird weiter in Abschnitt~\ref{sec:sekundaerliteratur} erläutert.
\end{itemize}


\section{Eintragstypen für die Literatur-Datenbank}

Die verwendete Literatur pflegen Sie in einer Literatur-Datenbank im Bibtex-Format. Dabei handelt es sich um eine Textdatei, wobei für jede Quelle mittels Name-Value-Pairs die relevanten Attribute (Autor, Titel etc.) hinterlegt sind. Die Datei wird üblicherweise nicht im Texteditor, sondern in einem spezialisierten Programm wie JabRef bearbeitet.

Sofern in der Literatur-Datenbank der Typ eines Eintrags (Entry Type) korrekt festgelegt ist, wird er im Literaturverzeichnis automatisch richtig dargestellt. Mit folgenden Typen sollten Sie i.d.R.\ auskommen:
\begin{description}
\item[article] Artikel in einer Fachzeitschrift, auch E-Journal (Zeitschrift in elektronischer Form)\footnote{Bei E-Journals/E-Books werden beim Zitieren anstelle der (u.U. nicht eindeutigen, da von der Schriftgröße abhängigen) Seitenzahl Abschnitt und Absatz näher bezeichnet, also: \cite[Abschnitt 1.2.3, Absatz 4]{Staab}.}
\item[book] Buch, auch E-Book 
\item[inbook] Kapitel in einem Buch, zu dem mehrere Autoren beigetragen haben 
\item[inproceedings] Beitrag zu einer Fachtagung/Konferenz 
\item[manual] Handbuch
\item[misc] anderweitig nicht zuordenbarer Typ
\item[phdthesis] Dissertation
\item[thesis] Bachelor-/Master-/Diplomarbeit (Art wird im Attribut \enquote{type} festgelegt) 
\item[online] Internet- oder Intranet-Quelle\footnote{\label{fn:onlineEntryType}Man beachte, dass der Eintragstyp \enquote{online} in JabRef nur im \enquote{biblatex-Modus} (Menü: Datei -- Neue biblatex Bibliothek) auswählbar ist.}
\item[report] technischer Bericht, Forschungsbericht oder White Paper; diesen Typ können Sie auch verwenden, um eine Projektarbeit zu zitieren (Art wird im Attribut \enquote{type} festgelegt) 
\end{description} 

Eine Übersicht über die notwendigen Attribute jedes Eintragstyps gibt die folgende Tabelle, wobei ein Schrägstrich als \enquote{oder} zu verstehen ist.\footnote{Auszugsweise entnommen aus \cite{biblatex:manual}.} Zudem sind die wichtigsten optionalen Attribute aufgeführt.

\begin{tabular}{l|p{6cm}|p{6cm}}
\textbf{Eintragstyp} & \textbf{notwendige Attribute} & \textbf{optionale Attribute (Auswahl)} \\
\hline
article & author, title, journal, year/date  & volume, number, pages, month, note \\
\hline
book & author, title, year/date  & publisher, edition, editor, \mbox{volume/number}, series, isbn, url \\
\hline
inbook & author, title, booktitle, year/date  & bookauthor, editor, volume/number, series, isbn, url \\
\hline
inproceedings & author, title, booktitle, year/date  & organization/publisher, editor, volume/number, series, isbn, url  \\
\hline
manual & author/editor, title, year/date  & organization/publisher, address, edition, month, note, url, urldate \\
\hline
misc & author/editor, title, year/date & howpublished, organization, month, note \\
\hline
phdthesis & author, title, institution, year/date  & address, month, note \\
\hline
thesis & author, title, institution, type, year/date &  address, month, note \\
\hline
online\textsuperscript{\ref{fn:onlineEntryType}} & author/editor, title, year\footnotemark/date, url & urldate \\
\hline
report & author, title, institution, type, year/date & number, version, url, urldate \\
\end{tabular}
\footnotetext{%
Sofern kein Jahr bekannt ist, sollte das Attribut nicht leer gelassen werden (sonst wird die aktuelle Jahreszahl automatisch eingefügt), sondern der Eintrag \enquote{o.J.} gewählt werden.}



\section{Zitieren von Sekundärliteratur}\label{sec:sekundaerliteratur}
Gelegentlich lässt es sich nicht vermeiden, aus der Sekundärliteratur zu zitieren. Dies leistet der folgende Befehl.
\begin{verbatim}
\footcitePrimaerSekundaer{Primaerquelle}{Seite}{Sekundaerquelle}{Seite}
\end{verbatim}
Die erste Seitenangabe bezieht sich auf die Primär-, die zweite auf die Sekundärquelle.
Die Seitenangaben sind optional, sie können auch leer bleiben.\footcitePrimaerSekundaer{Primaerquelle}{23}{Sekundaerquelle}{}
Es ist aber zu beachten, dass der Befehl \verb|\footcitePrimaerSekundaer| vier Argumente hat.

Ins Literaturverzeichnis soll nur die Sekundärquelle aufgenommen werden. Dies wird dadurch erreicht, dass in der Literatur-Datenbank bei der Primärquelle im Attribut \enquote{keyword} der Wert \enquote{ausblenden} eintragen wird.

\chapter{Beispiele für Abbildungen und Tabellen}\label{chapter:abbildungenTabellen}

Hier finden Sie Beispiele für Abbildungen, Tabellen, Formelsatz und  Source Code.

\section{Abbildungen}
In diesem Abschnitt gibt die Abbildungen~\ref{abb:Logo2cmHoch} und~\ref{abb:Logo2cmBreit}, die beide das Logo der DHBW zeigen.

\begin{figure}[htb]
\centering
\includegraphics[height=2cm]{graphics/dhbw.png}
\caption[DHBW-Logo 2cm hoch]{DHBW-Logo 2cm hoch.\footnotemark}
\label{abb:Logo2cmHoch}
\end{figure}
\footnotetext{Mit Änderungen entnommen aus: \cite{OhneAutorenOhneJahr}}

\lstset{language=TeX, % hervorzuhebende Keywords definieren
  morekeywords={footnotetext,footnotemark,footcite,caption}
}

\emph{Spezialfall:} Sofern \emph{innerhalb} der Bezeichnung einer Abbildung eine Fußnote angegeben oder eine Quelle referenziert werden soll, geschieht dies nicht per \lstinline|\footnote| oder \lstinline
|\footcite|. Vielmehr sind die Befehle \lstinline|\footnotemark| und \lstinline|\footnotetext| zu verwenden und außerdem das optionale Argument für \lstinline|\caption| anzugeben (vgl.\ Source Code).

\begin{figure}[htb]
\centering
\includegraphics[width=2cm]{graphics/dhbw.png}
\caption[DHBW-Logo 2cm breit.]{DHBW-Logo 2cm breit. (Quelle: DHBW\footnotemark)}
\label{abb:Logo2cmBreit}
\end{figure}
\footnotetext{\url{www.dhbw.de}}



\section{Tabellen}

In diesem Abschnitt gibt es zwei Beispiel-Tabellen, nämlich auf Seite~\pageref{tab:BeispielTabelleKlein} und auf Seite~\pageref{tab:BeispielTabelleGroesser}.

\begin{table}[htb]
\centering
\begin{tabular}{lcr}
links & Mitte & rechts \\
\hline
Muster & Muster & Muster \\
\end{tabular}
\caption{Kleine Beispiel-Tabelle.}
\label{tab:BeispielTabelleKlein}
\end{table}

\begin{table}[htb]
\centering
\begin{tabular}{|l|l|c|l|r||l}
    \textbf{Spalte 1} & \textbf{Spalte 2} & \textbf{Spalte 3} & \textbf{Spalte 4} & \textbf{Spalte 5} & \textbf{Spalte 6} \\
    \hline
    a        & b          & c                & d        & e        & f        \\
    Test     & Test, Test & Test, Test, Test & ~        & ~        & ~        \\
    1        & 2          & 3                & 4        & 5        & 6        \\
\end{tabular}
\caption{Größere Beispiel-Tabelle.}
\label{tab:BeispielTabelleGroesser}
\end{table}

\section{Etwas Mathematik}

Eine abgesetzte Formel:
\[
  \int_a^b x^2 \: \mathrm{d} x = \frac{1}{3} (b^3 - a^3)
\]

Es ist $a^2+b^2 = c^2$ eine Formel im Text.

\section{Source Code}

Source Code-Blöcke können auf folgende Arten eingefügt werden:

\lstset{language=Java}

Direkt im \LaTeX-Source Code:
\begin{lstlisting}
if(1 > 0) {
  System.out.println("OK"); 
} else {
  System.out.println("merkwuerdig");
}
\end{lstlisting}

oder eingefügt aus einer externen Datei.
\lstinputlisting{includes/HelloWorld.java}

\chapter*{Anhang}
\addcontentsline{toc}{chapter}{Anhang}
\section*{Anhangverzeichnis}
\vspace{-8em}

% vor \listofanhang müssen Einrückungen angepasst werden
\abstaendeanhangverzeichnis

\listofanhang
\clearpage
\spezialkopfzeile{Anhang} % damit in der Kopfzeile das Wort "Anhang" angezeigt wird

\anhang{So funktioniert's}

\lstset{language=TeX, % hervorzuhebende Keywords definieren
  morekeywords={anhang, anhangteil}
}


Um den Anforderungen der Zitierrichtlinien nachzukommen, wird das Paket \verb|tocloft| verwendet. Jeder Anhang wird mit dem (neu definierten) Befehl \lstinline|\anhang{Bezeichnung}| begonnen, der insbesondere dafür sorgt, dass ein Eintrag im Anhangsverzeichnis erzeugt wird. Manchmal ist es wünschenswert, auch einen Anhang noch weiter zu unterteilen. Hierfür wurde der Befehl \lstinline|\anhangteil{Bezeichnung}| definiert.

In~\ref{anhang:abbildung} finden Sie eine bekannte Abbildung und etwas Source Code in~\ref{anhang:sourcecode}. 

\anhangteil{Wieder mal eine Abbildung}\label{anhang:abbildung}
\begin{figure}[htb]
\centering
\includegraphics[width=0.9\linewidth]{graphics/dhbw.png}
\caption{Mal wieder das DHBW-Logo.}
\end{figure}

\anhangteil{Etwas Source Code}\label{anhang:sourcecode}
% \lstinputlisting{includes/HelloWorld.java}


\anhang{Release Notes}
\anhangteil{Änderungen in Version 1.1}\label{anhang:ReleaseNotes11}
In Version 1.1 sind einige Rückmeldungen, die nach der Einführungsvorlesung am 6.2.2015 oder nach Veröffentlichung der Vorlage in Moodle eingegangen sind, berücksichtigt worden. Korrekturen sind mit \enquote{(Fix)} gekennzeichnet. 

\begin{itemize}
\item \verb|latex-vorlage.tex|
\begin{itemize}
\item (Fix) Abkürzungsverzeichnis wird vor Abbildungsverzeichnis platziert
\item (Fix) Abbildungs- und Tabellenverzeichnis in Inhaltsverzeichnis aufgenommen
\item (Fix) Quellenverzeichnis wird nun ohne Kapitelnummer dargestellt

\item eingebundene Dateien in Unterverzeichnissen \verb|includes| bzw.\ \verb|graphics|
\item Beispiel-Anhang (Datei \verb|anhang.tex|) mit Erklärungen wurde eingebunden 
\end{itemize}

\item \verb|_dhbw_praeambel.tex|
\begin{itemize}
\item (Fix) das Paket hyperref wird nach biblatex eingebunden, um ein Problem mit der Verlinkung der Fußnoten im PDF zu beheben
\item (Fix) Fußnoten  gemäß der Richtlinien fortlaufend nummeriert und nicht pro Kapitel
\item Einstellungen hinzugefügt, um Anhangsverzeichnis zu ermöglichen
\item bessere Kompatibilität zwischen KOMA-Script (scrreprt) und anderen Paketen mittels scrhack
\end{itemize}

\item \verb|_dhbw_biblatex-config.tex|
\begin{itemize}
\item (Fix) keine Abschnittsnummern für einzelne Verzeichnisse im Quellenverzeichnis
\end{itemize}

\item \verb|abbildungen_und_tabellen.tex|
\begin{itemize}
\item Erklärung, wie eine Fußnote/ein Zitat bei einer Abbildung zu erstellen ist
\end{itemize}

\item \verb|abkuerzungen.tex|
\begin{itemize}
\item Abkürzungsverzeichnis wird im Inhaltsverzeichnis aufgeführt
\end{itemize}

\item \verb|abstract.tex|, \verb|anhang.tex|, \verb|einleitung.tex| 
\begin{itemize}
\item Erklärungen im Text ergänzt
\end{itemize}

\item \verb|deckblatt.tex|
\begin{itemize}
\item Meta-Daten (Autor, Titel) für die generierte PDF-Datei lassen sich nun festlegen
\end{itemize}

\end{itemize}


\anhangteil{Änderungen in Version 1.2}\label{anhang:ReleaseNotes12}
Über das Forum in Moodle sind einige Rückmeldungen eingegangen -- vielen Dank an alle, die dazu beigetragen haben. In der Version 1.2 wurden folgende Änderungen vorgenommen, wobei Korrekturen wieder mit \enquote{(Fix)} gekennzeichnet sind: 

\begin{itemize}
\item \verb|latex-vorlage.tex| (Hauptdokument)
\begin{itemize}
\item (Fix) Zeile 19: Seitenzahlen zu Beginn mit römischen \emph{Groß}buchstaben nummeriert
\end{itemize}

\item \verb|_dhbw_praeambel.tex|
\begin{itemize}
\item Zeile 39/40: Unterstützung für \enquote{ebenda} 
\item Zeile 46--68: zweite Gliederungsebene für Anhänge ermöglicht
\item (Fix) Zeile 70--73: Abbildungen und Tabellen: Zähler fortlaufend, kein Rücksetzen zu Kapitelbeginn (Paket \verb|chngcntr| anstelle von Paket \verb|remreset|)
\end{itemize}

\item \verb|_dhbw_biblatex-config.tex|
\begin{itemize}
\item (Fix) bei Quellen mit Herausgeber, aber ohne Autor wird der Name des Herausgebers im Verzeichnis fett gedruckt
\item Unterstützung für \enquote{ebenda} 
\end{itemize}

\item \verb|abkuerzungen.tex|
\begin{itemize}
\item Bemerkungen zur fortgeschrittenen Nutzung des \verb|acronym|-Pakets eingefügt 
\end{itemize}

\item \verb|einleitung.tex|
\begin{itemize}
\item Abschnitt 1.3 zu Einstellungen ergänzt
\item Abschnitt 1.5 zu Fehlerbehebungen eingefügt 
\end{itemize}

\item \verb|text-mit-zitaten.tex|
\begin{itemize}
\item Abschnitt 3.1 eingefügt, Erläuterungen zum Zitieren mit \enquote{vgl.} und \enquote{ebenda}. 
\item Abschnitt 3.2: Beispiele ergänzt
\item Hinweis zu Jahreszahlen bei Online-Quellen
\end{itemize}

\item \verb|anhang.tex|
\begin{itemize}
\item Erläuterungen zur zweiten Gliederungsebene
\end{itemize}

\item \verb|literatur-datenbank.bib|
\begin{itemize}
\item weitere Beispiele für Quellen
\end{itemize}

\end{itemize}

\anhangteil{Änderungen in Version 1.3}\label{anhang:ReleaseNotes13}
Durch die ab 1/2016 geltenden Änderungen der Zitierrichtlinien des Studiengangs waren einige kleinere Anpassungen der Vorlage erforderlich, die nachfolgend beschrieben sind. Bei dieser Gelegenheit ebenfalls erfolgte Korrekturen sind wieder mit \enquote{(Fix)} gekennzeichnet:

\begin{itemize}
\item \verb|latex-vorlage.tex| (Hauptdokument)
\begin{itemize}
\item Hinweis auf Option doppelseitiger Druck entfernt
\item Schriftgröße der Kapitelüberschriften verkleinert
\item (Fix) Kopf- und Fußzeilen werden nun korrekt angezeigt für erste Seite eines Kapitels und auch  Quellenverzeichnisse
\end{itemize}

\item \verb|_dhbw_praeambel.tex|
\begin{itemize}
\item Angabe des unteren Rands für Seitenzahl, da diese nun unten rechts steht
\item Unterstützung für \enquote{ebenda} entfernt
\item (Fix) Präfixe wie \enquote{von} im Namen eines Autors werden berücksichtigt
\item Anpassung der Abstände bei Kapitelüberschriften
\item Kopf- und Fußzeile für Verzeichnisse nun in \verb|_dhbw_kopfzeilen.tex| definiert 
\end{itemize}


\item \verb|deckblatt.tex|
\begin{itemize}
\item Schriftgröße des Titels vergrößert
\item Befehl \verb|\typMeinerArbeit| eingeführt, um Typ auszuwählen
\item Festlegung des Themas (für ehrenwörtliche Erklärung) mit Befehl \verb|\themaMeinerArbeit|
\item Darstellung der Angabe des Betreuers in der Ausbildungsstätte angepasst
\item Formulierung des Sperrvermerks angepasst  
\end{itemize}

\item \verb|_dhbw_erklaerung.tex|
\begin{itemize}
\item Formulierung angepasst an geänderte Prüfungsordnung
\item Typ und Thema der Arbeit werden automatisch eingefügt
\end{itemize}

\item \verb|_dhbw_kopfzeilen.tex|
\begin{itemize}
\item Seitennummern stehen jetzt unten rechts
\item (Fix) Kopf- und Fußzeile werden nun korrekt angezeigt in Verzeichnissen und dem Anhang
\end{itemize}

\item \verb|_dhbw_biblatex-config.tex|
\begin{itemize}
\item Anpassung des Zitierstils auf die ab 1/2016 geltenden Regelungen  
\item Vorkehrungen für Eindeutigkeit (Hinzufügen abgekürzter oder nötigenfalls ausgeschriebener Vorname) bei Übereinstimmung von Name und Jahreszahl 
\end{itemize}

\item \verb|einleitung.tex|
\begin{itemize}
\item Abschnitt 1.3 zu Einstellungen grundlegend überarbeitet
\item Abschnitt 1.5.2 zur Kontrolle der Seitenränder eingefügt 
\end{itemize}

\item \verb|text-mit-zitaten.tex|
\begin{itemize}
\item Abschnitt 3.1: Hinweise zu \enquote{ebenda} entfernt
\item Abschnitt 3.2: Beispiele zur Eindeutigkeit des Zitats ergänzt
\item Abschnitt 3.3: Hinweise für E-Journals/E-Books ergänzt 
\end{itemize}

\item \verb|anhang.tex|
\begin{itemize}
\item (Fix) Befehl \verb|\spezialkopfzeile| aufgenommen, damit in Kopfzeile das Wort \enquote{Anhang} angezeigt wird 
\item diese Release Notes wurden in eine eigene Datei verschoben
\end{itemize}

\item \verb|release_notes.tex|
\begin{itemize}
\item s.o.
\end{itemize}


\item \verb|literatur-datenbank.bib|
\begin{itemize}
\item weitere Beispiele für Quellen
\end{itemize}
\end{itemize}

\anhangteil{Änderungen in Version 1.4}\label{anhang:ReleaseNotes14}
Durch nicht abwärtskompatible Änderungen beim Versionswechsel von Biblatex 3.2 zu 3.3 sind einige Änderungen notwendig geworden.\footnote{Diese basieren auf Vorschlägen von Yannik Ehlert -- vielen Dank dafür!}
Die vorliegende Version 1.4 wurde erfolgreich mit MikTeX gestestet (portable Version 2.9.6361 vom 3.6.2017, unter Verwendung von Biblatex 3.7).

\begin{itemize}
\item \verb|_dhbw_biblatex-config.tex|
\begin{itemize}
\item Anpassung der \verb|\usebibmacro|-Befehle
\end{itemize}

\item \verb|_dhbw_authoryear.bbx|
\begin{itemize}
\item  Änderung von \verb|\printdateextralabel| zu \verb|\printlabeldateextra|
\end{itemize}
\end{itemize}

\anhangteil{Änderungen in Version 1.5}\label{anhang:ReleaseNotes15}
Für den Test dieser Version auf einem Windows-System wurde wieder die portable Version von MiKTeX (2.9.6521 vom 10.11.2017) verwendet.\footnote{\url{http://miktex.org/portable}} Da in diesem Paket leider die Versionen von Biblatex (3.10) und Biber (2.7) inkompatibel sind, ist es erforderlich, die Datei \verb|biber.exe| im Verzeichnis \verb|texmfs\install\miktex\bin\| durch die aktuelle Version 2.10 vom 20.12.2017\footnote{\url{https://sourceforge.net/projects/biblatex-biber/files/biblatex-biber/current/binaries/Windows/}} zu ersetzen. Im Editor TeXworks verwendet man dann zum Übersetzen des \LaTeX-Sourcecodes Typeset/pdfLaTeX bzw.\ Typeset/Biber.

Korrekturen sind wieder mit \enquote{(Fix)} gekennzeichnet.

\begin{itemize}
\item \verb|latex-vorlage.tex| (Hauptdokument)
\begin{itemize}
\item Nach der Änderung der Zitierrichtlinien gibt es nun kein separates Verzeichnis mehr für Internet- und Intranetquellen.
\item Option \verb|notkeyword=ausblenden| bei \verb|\printbibligraphy| sorgt dafür, dass Sekundärliteratur korrekt zitiert wird.
\end{itemize}

\item \verb|_dhbw_praembel.tex|
\begin{itemize}
\item (Fix) Die Bezeichnung geschachtelter Anhänge wurde auf das in den Zitierrichtlinien geforderte Format \enquote{Anhang 2/1} angepasst (Befehl \verb|\anhangteil|).
\end{itemize}

\item \verb|einleitung.tex|
\begin{itemize}
\item Hinweis zum Ausblenden der farbigen Links im PDF hinzugefügt
\end{itemize}

\item \verb|text-mit-zitaten.tex|
\begin{itemize}
\item Abschnitt 3.4 aktualisiert nach Wegfall des separaten Verzeichnisses für Internet- und Intranetquellen
\item Abschnitt zum Zitieren von Sekundärliteratur hinzugefügt
\end{itemize}

\end{itemize}


\anhangteil{Änderungen in Version 1.6}\label{anhang:ReleaseNotes16}
Diese Version wurde auf einem Windows-System erfolgreich mit der portablen Version von MiKTeX (2.9.6621 vom 18.02.2018) getestet.\footnote{Vielen Dank an Florian Eichin für seine wertvollen Anmerkungen.}

Korrekturen sind wieder mit \enquote{(Fix)} gekennzeichnet.

\newpage

\begin{itemize}
\item \verb|latex-vorlage.tex| (Hauptdokument)
\begin{itemize}
\item (Fix) An einer Stelle gab es in Version 1.5 (Internetquellen nicht mehr separat) noch ein Überbleibsel von Version 1.4 (Internetquellen separat), dies wurde korrigiert.
\item (Fix) Im Inhaltsverzeichnis war die Verlinkung des Abbildungs- und Tabellenverzeich\-nisses nicht ganz korrekt.
\item Mit den Befehlen \verb|\literaturverzeichnis| bzw.\ \verb|\literaturUndQuellenverzeichnis| kann bequem die Erstellung der Quellenverzeichnisse gesteuert werden, abhängig davon, ob es ein Gesprächsverzeichnis gibt oder nicht.
 
\end{itemize}

\item \verb|_dhbw_praembel.tex|
\begin{itemize}
\item Einrückungen für Abbildungs-, Tabellen- und Anhangverzeichnis angepasst
\item Abkürzungen \enquote{Abb.} und \enquote{Tab.} für Abbildungen bzw.\ Tabellen
\end{itemize}

\item \verb|_dhbw_biblatex-config.tex|
\begin{itemize}
\item Befehle \verb|\literaturverzeichnis| und \verb|\literaturUndGespraechsverzeichnis| definiert
\item Befehl \verb|\footcitePrimaerSekundaer| definiert
\end{itemize}

\item \verb|_dhbw_erklaerung.tex|
\begin{itemize}
\item Eintrag als \enquote{Erklärung} (statt \enquote{Ehrenwörtliche Erklärung}) ins Inhaltsverzeichnis
\end{itemize}

\item \verb|einleitung.tex|
\begin{itemize}
\item Bezeichnung \enquote{Erklärung} statt \enquote{Ehrenwörtliche Erklärung}
\item Erläuterung von \verb|\literaturverzeichnis| und \verb|\literaturUndGespraechsverzeichnis|
\item Hinweis auf Notwendigkeit von Updates bei MikTeX Portable
\end{itemize}

\item \verb|text_mit_zitaten.tex|
\begin{itemize}
\item Erläuterungen zu Befehl \verb|\footcitePrimaerSekundaer| ergänzt
\end{itemize}

\item \verb|anhang.tex|
\begin{itemize}
\item Befehl \verb|\abstaendeanhangverzeichnis| für Anpassung Einrückung ergänzt
\end{itemize}

\item \verb|literatur-datenbank.bib|
\begin{itemize}
\item Eintrag ergänzt
\end{itemize}

\end{itemize}

\anhangteil{Änderungen in Version 1.7}\label{anhang:ReleaseNotes17}
Diese Version wurde auf einem Windows-System erfolgreich mit der portablen Version von MiKTeX (2.9.6942 vom 04.01.2019) getestet.

Korrekturen sind wieder mit \enquote{(Fix)} gekennzeichnet.

\begin{itemize}
\item \verb|_dhbw-authoryear.bbx|
\begin{itemize}
\item Da \verb|labeldate| in Biblatex nicht mehr unterstützt wird, erfolgte eine Umbenennung in 
\verb|labeldateparts|.\footnote{vgl.\ \url{https://github.com/semprag/biblatex-sp-unified/issues/23}}
\end{itemize}

\item \verb|_dhbw_biblatex-config.tex|
\begin{itemize}
\item (Fix) Es wurde das Problem behoben, dass im Literaturverzeichnis bei bestimmten Eintragstypen der Titel in Anführungszeichen steht.\footnote{Danke an Florian Eichin für seinen Hinweis.}
\end{itemize}

\end{itemize}


\anhangteil{Änderungen in Version 1.8}\label{anhang:ReleaseNotes18}
Diese Version wurde auf einem Windows-System erfolgreich mit der portablen Version von MiKTeX (2.9.6942 vom 04.01.2019) getestet.

Die Aktualisierungen in der Vorlage spiegeln zum Einen die Änderungen in den Zitierrichtlinien wieder. Zum Anderen wurden einige studentische Vorschläge aufgegriffen, um die Nutzung der Vorlage zu erleichtern.\footnote{Danke an Bjarne Koll, Tobias Schwarz und Lars Ungerathen für ihre Anregungen.} 

\begin{itemize}

\item \verb|latex_vorlage.tex| (Hauptdokument)
\begin{itemize}
\item Es wird nun davon ausgegangen, dass die zur Vorlage gehörenden Dateien in einem eigenen Verzeichnis (\verb|template|) liegen.
\item Stellenweise wurden Erläuterungen als Kommentare hinzugefügt.
\end{itemize}

\item \verb|_dhbw_biblatex-config.tex|
\begin{itemize}
\item Code, der mehrere Quellenverzeichnisse unterstützt, wurde entfernt.
\item Ein zu großer Abstand nach Zitaten von Sekundärliteratur wurde korrigiert. 
\end{itemize}

\item \verb|_dhbw_erklaerung.bbx|
\begin{itemize}
\item Gemäß der Anforderung in den Zitierrichtlinien wird die Erklärung nicht ins Inhaltsverzeichnis aufgenommen und nicht mit einer Seitenzahl versehen. 
\end{itemize}

\pagebreak
\item \verb|_dhbw_praeambel.bbx|
\begin{itemize}
\item Gemäß der Anforderung in den Zitierrichtlinien werden im Literaturverzeichnis alle Autor/innen eines Werks angegeben.
\end{itemize}

\item \verb|abstract.tex|
\begin{itemize}
\item Hinweis auf \LaTeX-Spickzettel hinzugefügt.
\end{itemize}

\item \verb|deckblatt.tex|
\begin{itemize}
\item Vorname, Name, Titel der Arbeit sind nur zu Beginn einzutragen und werden dann an den entsprechenden Stellen automatisch ergänzt.
\item Hervorhebung, dass Angaben zum Unternehmen sowie den Betreuer/innen zu ergänzen sind. 
\item Wortlaut des Vertraulichkeitsvermerks wurde an die aktuelle Fassung in der Studien- und Prüfungsordnung angepasst. 
\end{itemize}

\item \verb|einleitung.tex|
\begin{itemize}
\item Ein eigenständiges Gesprächsverzeichnis als Teil des Quellenverzeichnisses ist in den Zitierrichtlinien nicht mehr vorgesehen, die entsprechenden Hinweise wurden entfernt.
\item Ein alter Hinweis auf die Darstellung von Links im Verzeichnis der Internetquellen wurde entfernt, da es ein solches eigenständiges Verzeichnis nicht mehr gibt. 
\end{itemize}

\item \verb|text_mit_zitaten.tex|
\begin{itemize}
\item Es wird nun erläutert, wie zwei Quellenangaben unmittelbar nebeneinander dargestellt werden können.
\item Erklärungen, die von mehreren Quellenverzeichnissen ausgegangen sind, wurden entfernt.
\end{itemize}

\item \verb|literatur-datenbank.bib|
\begin{itemize}
\item Gespräch wurde entfernt, da dieses nicht mehr im Quellenverzeichnis aufgeführt werden soll.
\end{itemize}

\end{itemize}
%%% Ende des eigentlichen Inhalts %%%


%%% Quellenverzeichnisse (keine Anpassung nötig) %%%
\clearpage
\literaturverzeichnis
%%% Ende Quellenverzeichnisse %%%


%%% Erklärung (keine Anpassungen nötig) %%%
% steht ganz am Ende des Dokuments
\cleardoublepage
\clearpage

\thispagestyle{empty}

{\LARGE\textsf{\textbf{Erklärung}}\bigskip}

% \typMeinerArbeit und \themaMeinerArbeit werden in deckblatt.tex definiert
Ich versichere hiermit, dass ich meine \typMeinerArbeit\ mit dem Thema: \emph{\themaMeinerArbeit} selbstständig verfasst und keine anderen als die angegebenen Quellen und Hilfsmittel benutzt habe.
Ich versichere zudem, dass die eingereichte elektronische Fassung mit der gedruckten Fassung übereinstimmt.

\vspace{3cm}

\begin{center}
\begin{tabular}{ccc}
(Ort, Datum) & \hspace{0.3\linewidth} & (Unterschrift)
\end{tabular}
\end{center}
\end{document}